\chapter{Model Providers}

% === Source file: providers/anthropic.md ===
\section{Anthropic (Claude)}


Anthropic builds the \textbf{Claude} model family and provides access via an API.
In OpenClaw you can authenticate with an API key or a \textbf{setup-token}.

\subsection{Option A: Anthropic API key}


\textbf{Best for:} standard API access and usage-based billing.
Create your API key in the Anthropic Console.

\subsubsection{CLI setup}


\begin{lstlisting}[language=bash]
openclaw onboard
# choose: Anthropic API key

# or non-interactive
openclaw onboard --anthropic-api-key "$ANTHROPIC_API_KEY"
\end{lstlisting}


\subsubsection{Config snippet}


\begin{lstlisting}[language=json]
{
  env: { ANTHROPIC_API_KEY: "sk-ant-..." },
  agents: { defaults: { model: { primary: "anthropic/claude-opus-4-6" } } },
}
\end{lstlisting}


\subsection{Thinking defaults (Claude 4.6)}


\begin{itemize}
  \item Anthropic Claude 4.6 models default to \texttt{adaptive} thinking in OpenClaw when no explicit thinking level is set.
  \item You can override per-message (\texttt{/think:}) or in model params:
\end{itemize}

\texttt{agents.defaults.models["anthropic/"].params.thinking}.
\begin{itemize}
  \item Related Anthropic docs:
  \item \href{https://platform.claude.com/docs/en/build-with-claude/adaptive-thinking}{Adaptive thinking}
  \item \href{https://platform.claude.com/docs/en/build-with-claude/extended-thinking}{Extended thinking}
\end{itemize}


\subsection{Prompt caching (Anthropic API)}


OpenClaw supports Anthropic's prompt caching feature. This is \textbf{API-only}; subscription auth does not honor cache settings.

\subsubsection{Configuration}


Use the \texttt{cacheRetention} parameter in your model config:


\begin{table}[h]
\centering
\small
\begin{tabular}{|l|l|l|}
\hline
Value & Cache Duration & Description \\
\hline
\texttt{none} & No caching & Disable prompt caching \\
\hline
\texttt{short} & 5 minutes & Default for API Key auth \\
\hline
\texttt{long} & 1 hour & Extended cache (requires beta flag) \\
\hline
\end{tabular}
\end{table}



\begin{lstlisting}[language=json]
{
  agents: {
    defaults: {
      models: {
        "anthropic/claude-opus-4-6": {
          params: { cacheRetention: "long" },
        },
      },
    },
  },
}
\end{lstlisting}


\subsubsection{Defaults}


When using Anthropic API Key authentication, OpenClaw automatically applies \texttt{cacheRetention: "short"} (5-minute cache) for all Anthropic models. You can override this by explicitly setting \texttt{cacheRetention} in your config.

\subsubsection{Per-agent cacheRetention overrides}


Use model-level params as your baseline, then override specific agents via \texttt{agents.list[].params}.

\begin{lstlisting}[language=json]
{
  agents: {
    defaults: {
      model: { primary: "anthropic/claude-opus-4-6" },
      models: {
        "anthropic/claude-opus-4-6": {
          params: { cacheRetention: "long" }, // baseline for most agents
        },
      },
    },
    list: [
      { id: "research", default: true },
      { id: "alerts", params: { cacheRetention: "none" } }, // override for this agent only
    ],
  },
}
\end{lstlisting}


Config merge order for cache-related params:

\begin{enumerate}
  \item \texttt{agents.defaults.models["provider/model"].params}
  \item \texttt{agents.list[].params} (matching \texttt{id}, overrides by key)
\end{enumerate}


This lets one agent keep a long-lived cache while another agent on the same model disables caching to avoid write costs on bursty/low-reuse traffic.

\subsubsection{Bedrock Claude notes}


\begin{itemize}
  \item Anthropic Claude models on Bedrock (\texttt{amazon-bedrock/\textit{anthropic.claude}}) accept \texttt{cacheRetention} pass-through when configured.
  \item Non-Anthropic Bedrock models are forced to \texttt{cacheRetention: "none"} at runtime.
  \item Anthropic API-key smart defaults also seed \texttt{cacheRetention: "short"} for Claude-on-Bedrock model refs when no explicit value is set.
\end{itemize}


\subsubsection{Legacy parameter}


The older \texttt{cacheControlTtl} parameter is still supported for backwards compatibility:

\begin{itemize}
  \item \texttt{"5m"} maps to \texttt{short}
  \item \texttt{"1h"} maps to \texttt{long}
\end{itemize}


We recommend migrating to the new \texttt{cacheRetention} parameter.

OpenClaw includes the \texttt{extended-cache-ttl-2025-04-11} beta flag for Anthropic API
requests; keep it if you override provider headers (see \href{/gateway/configuration}{/gateway/configuration}).

\subsection{1M context window (Anthropic beta)}


Anthropic's 1M context window is beta-gated. In OpenClaw, enable it per model
with \texttt{params.context1m: true} for supported Opus/Sonnet models.

\begin{lstlisting}[language=json]
{
  agents: {
    defaults: {
      models: {
        "anthropic/claude-opus-4-6": {
          params: { context1m: true },
        },
      },
    },
  },
}
\end{lstlisting}


OpenClaw maps this to \texttt{anthropic-beta: context-1m-2025-08-07} on Anthropic
requests.

This only activates when \texttt{params.context1m} is explicitly set to \texttt{true} for
that model.

Requirement: Anthropic must allow long-context usage on that credential
(typically API key billing, or a subscription account with Extra Usage
enabled). Otherwise Anthropic returns:
\texttt{HTTP 429: rate\_limit\_error: Extra usage is required for long context requests}.

Note: Anthropic currently rejects \texttt{context-1m-*} beta requests when using
OAuth/subscription tokens (\texttt{sk-ant-oat-*}). OpenClaw automatically skips the
context1m beta header for OAuth auth and keeps the required OAuth betas.

\subsection{Option B: Claude setup-token}


\textbf{Best for:} using your Claude subscription.

\subsubsection{Where to get a setup-token}


Setup-tokens are created by the \textbf{Claude Code CLI}, not the Anthropic Console. You can run this on \textbf{any machine}:

\begin{lstlisting}[language=bash]
claude setup-token
\end{lstlisting}


Paste the token into OpenClaw (wizard: \textbf{Anthropic token (paste setup-token)}), or run it on the gateway host:

\begin{lstlisting}[language=bash]
openclaw models auth setup-token --provider anthropic
\end{lstlisting}


If you generated the token on a different machine, paste it:

\begin{lstlisting}[language=bash]
openclaw models auth paste-token --provider anthropic
\end{lstlisting}


\subsubsection{CLI setup (setup-token)}


\begin{lstlisting}[language=bash]
# Paste a setup-token during onboarding
openclaw onboard --auth-choice setup-token
\end{lstlisting}


\subsubsection{Config snippet (setup-token)}


\begin{lstlisting}[language=json]
{
  agents: { defaults: { model: { primary: "anthropic/claude-opus-4-6" } } },
}
\end{lstlisting}


\subsection{Notes}


\begin{itemize}
  \item Generate the setup-token with \texttt{claude setup-token} and paste it, or run \texttt{openclaw models auth setup-token} on the gateway host.
  \item If you see “OAuth token refresh failed …” on a Claude subscription, re-auth with a setup-token. See \href{/gateway/troubleshooting\#oauth-token-refresh-failed-anthropic-claude-subscription}{/gateway/troubleshooting\#oauth-token-refresh-failed-anthropic-claude-subscription}.
  \item Auth details + reuse rules are in \href{/concepts/oauth}{/concepts/oauth}.
\end{itemize}


\subsection{Troubleshooting}


\textbf{401 errors / token suddenly invalid}

\begin{itemize}
  \item Claude subscription auth can expire or be revoked. Re-run \texttt{claude setup-token}
\end{itemize}

and paste it into the \textbf{gateway host}.
\begin{itemize}
  \item If the Claude CLI login lives on a different machine, use
\end{itemize}

\texttt{openclaw models auth paste-token --provider anthropic} on the gateway host.

\textbf{No API key found for provider "anthropic"}

\begin{itemize}
  \item Auth is \textbf{per agent}. New agents don’t inherit the main agent’s keys.
  \item Re-run onboarding for that agent, or paste a setup-token / API key on the
\end{itemize}

gateway host, then verify with \texttt{openclaw models status}.

\textbf{No credentials found for profile \texttt{anthropic:default}}

\begin{itemize}
  \item Run \texttt{openclaw models status} to see which auth profile is active.
  \item Re-run onboarding, or paste a setup-token / API key for that profile.
\end{itemize}


\textbf{No available auth profile (all in cooldown/unavailable)}

\begin{itemize}
  \item Check \texttt{openclaw models status --json} for \texttt{auth.unusableProfiles}.
  \item Add another Anthropic profile or wait for cooldown.
\end{itemize}


More: \href{/gateway/troubleshooting}{/gateway/troubleshooting} and \href{/help/faq}{/help/faq}.


\clearpage

% === Source file: providers/bedrock.md ===
\section{Amazon Bedrock}


OpenClaw can use \textbf{Amazon Bedrock} models via pi‑ai’s \textbf{Bedrock Converse}
streaming provider. Bedrock auth uses the \textbf{AWS SDK default credential chain},
not an API key.

\subsection{What pi‑ai supports}


\begin{itemize}
  \item Provider: \texttt{amazon-bedrock}
  \item API: \texttt{bedrock-converse-stream}
  \item Auth: AWS credentials (env vars, shared config, or instance role)
  \item Region: \texttt{AWS\_REGION} or \texttt{AWS\_DEFAULT\_REGION} (default: \texttt{us-east-1})
\end{itemize}


\subsection{Automatic model discovery}


If AWS credentials are detected, OpenClaw can automatically discover Bedrock
models that support \textbf{streaming} and \textbf{text output}. Discovery uses
\texttt{bedrock:ListFoundationModels} and is cached (default: 1 hour).

Config options live under \texttt{models.bedrockDiscovery}:

\begin{lstlisting}[language=json]
{
  models: {
    bedrockDiscovery: {
      enabled: true,
      region: "us-east-1",
      providerFilter: ["anthropic", "amazon"],
      refreshInterval: 3600,
      defaultContextWindow: 32000,
      defaultMaxTokens: 4096,
    },
  },
}
\end{lstlisting}


Notes:

\begin{itemize}
  \item \texttt{enabled} defaults to \texttt{true} when AWS credentials are present.
  \item \texttt{region} defaults to \texttt{AWS\_REGION} or \texttt{AWS\_DEFAULT\_REGION}, then \texttt{us-east-1}.
  \item \texttt{providerFilter} matches Bedrock provider names (for example \texttt{anthropic}).
  \item \texttt{refreshInterval} is seconds; set to \texttt{0} to disable caching.
  \item \texttt{defaultContextWindow} (default: \texttt{32000}) and \texttt{defaultMaxTokens} (default: \texttt{4096})
\end{itemize}

are used for discovered models (override if you know your model limits).

\subsection{Onboarding}


\begin{enumerate}
  \item Ensure AWS credentials are available on the \textbf{gateway host}:
\end{enumerate}


\begin{lstlisting}[language=bash]
export AWS_ACCESS_KEY_ID="AKIA..."
export AWS_SECRET_ACCESS_KEY="..."
export AWS_REGION="us-east-1"
# Optional:
export AWS_SESSION_TOKEN="..."
export AWS_PROFILE="your-profile"
# Optional (Bedrock API key/bearer token):
export AWS_BEARER_TOKEN_BEDROCK="..."
\end{lstlisting}


\begin{enumerate}
  \item Add a Bedrock provider and model to your config (no \texttt{apiKey} required):
\end{enumerate}


\begin{lstlisting}[language=json]
{
  models: {
    providers: {
      "amazon-bedrock": {
        baseUrl: "https://bedrock-runtime.us-east-1.amazonaws.com",
        api: "bedrock-converse-stream",
        auth: "aws-sdk",
        models: [
          {
            id: "us.anthropic.claude-opus-4-6-v1:0",
            name: "Claude Opus 4.6 (Bedrock)",
            reasoning: true,
            input: ["text", "image"],
            cost: { input: 0, output: 0, cacheRead: 0, cacheWrite: 0 },
            contextWindow: 200000,
            maxTokens: 8192,
          },
        ],
      },
    },
  },
  agents: {
    defaults: {
      model: { primary: "amazon-bedrock/us.anthropic.claude-opus-4-6-v1:0" },
    },
  },
}
\end{lstlisting}


\subsection{EC2 Instance Roles}


When running OpenClaw on an EC2 instance with an IAM role attached, the AWS SDK
will automatically use the instance metadata service (IMDS) for authentication.
However, OpenClaw's credential detection currently only checks for environment
variables, not IMDS credentials.

\textbf{Workaround:} Set \texttt{AWS\_PROFILE=default} to signal that AWS credentials are
available. The actual authentication still uses the instance role via IMDS.

\begin{lstlisting}[language=bash]
# Add to ~/.bashrc or your shell profile
export AWS_PROFILE=default
export AWS_REGION=us-east-1
\end{lstlisting}


\textbf{Required IAM permissions} for the EC2 instance role:

\begin{itemize}
  \item \texttt{bedrock:InvokeModel}
  \item \texttt{bedrock:InvokeModelWithResponseStream}
  \item \texttt{bedrock:ListFoundationModels} (for automatic discovery)
\end{itemize}


Or attach the managed policy \texttt{AmazonBedrockFullAccess}.

\subsection{Quick setup (AWS path)}


\begin{lstlisting}[language=bash]
# 1. Create IAM role and instance profile
aws iam create-role --role-name EC2-Bedrock-Access \
  --assume-role-policy-document '{
    "Version": "2012-10-17",
    "Statement": [{
      "Effect": "Allow",
      "Principal": {"Service": "ec2.amazonaws.com"},
      "Action": "sts:AssumeRole"
    }]
  }'

aws iam attach-role-policy --role-name EC2-Bedrock-Access \
  --policy-arn arn:aws:iam::aws:policy/AmazonBedrockFullAccess

aws iam create-instance-profile --instance-profile-name EC2-Bedrock-Access
aws iam add-role-to-instance-profile \
  --instance-profile-name EC2-Bedrock-Access \
  --role-name EC2-Bedrock-Access

# 2. Attach to your EC2 instance
aws ec2 associate-iam-instance-profile \
  --instance-id i-xxxxx \
  --iam-instance-profile Name=EC2-Bedrock-Access

# 3. On the EC2 instance, enable discovery
openclaw config set models.bedrockDiscovery.enabled true
openclaw config set models.bedrockDiscovery.region us-east-1

# 4. Set the workaround env vars
echo 'export AWS_PROFILE=default' >> ~/.bashrc
echo 'export AWS_REGION=us-east-1' >> ~/.bashrc
source ~/.bashrc

# 5. Verify models are discovered
openclaw models list
\end{lstlisting}


\subsection{Notes}


\begin{itemize}
  \item Bedrock requires \textbf{model access} enabled in your AWS account/region.
  \item Automatic discovery needs the \texttt{bedrock:ListFoundationModels} permission.
  \item If you use profiles, set \texttt{AWS\_PROFILE} on the gateway host.
  \item OpenClaw surfaces the credential source in this order: \texttt{AWS\_BEARER\_TOKEN\_BEDROCK},
\end{itemize}

then \texttt{AWS\_ACCESS\_KEY\_ID} + \texttt{AWS\_SECRET\_ACCESS\_KEY}, then \texttt{AWS\_PROFILE}, then the
default AWS SDK chain.
\begin{itemize}
  \item Reasoning support depends on the model; check the Bedrock model card for
\end{itemize}

current capabilities.
\begin{itemize}
  \item If you prefer a managed key flow, you can also place an OpenAI‑compatible
\end{itemize}

proxy in front of Bedrock and configure it as an OpenAI provider instead.


\clearpage

% === Source file: providers/claude-max-api-proxy.md ===
\section{Claude Max API Proxy}


\textbf{claude-max-api-proxy} is a community tool that exposes your Claude Max/Pro subscription as an OpenAI-compatible API endpoint. This allows you to use your subscription with any tool that supports the OpenAI API format.

This path is technical compatibility only. Anthropic has blocked some subscription
usage outside Claude Code in the past. You must decide for yourself whether to use
it and verify Anthropic's current terms before relying on it.

\subsection{Why Use This?}



\begin{table}[h]
\centering
\small
\begin{tabular}{|l|l|l|}
\hline
Approach & Cost & Best For \\
\hline
Anthropic API & Pay per token (\textasciitilde{}\$15/M input, \$75/M output for Opus) & Production apps, high volume \\
\hline
Claude Max subscription & \$200/month flat & Personal use, development, unlimited usage \\
\hline
\end{tabular}
\end{table}



If you have a Claude Max subscription and want to use it with OpenAI-compatible tools, this proxy may reduce cost for some workflows. API keys remain the clearer policy path for production use.

\subsection{How It Works}


\begin{lstlisting}
Your App → claude-max-api-proxy → Claude Code CLI → Anthropic (via subscription)
     (OpenAI format)              (converts format)      (uses your login)
\end{lstlisting}


The proxy:

\begin{enumerate}
  \item Accepts OpenAI-format requests at \texttt{http://localhost:3456/v1/chat/completions}
  \item Converts them to Claude Code CLI commands
  \item Returns responses in OpenAI format (streaming supported)
\end{enumerate}


\subsection{Installation}


\begin{lstlisting}[language=bash]
# Requires Node.js 20+ and Claude Code CLI
npm install -g claude-max-api-proxy

# Verify Claude CLI is authenticated
claude --version
\end{lstlisting}


\subsection{Usage}


\subsubsection{Start the server}


\begin{lstlisting}[language=bash]
claude-max-api
# Server runs at http://localhost:3456
\end{lstlisting}


\subsubsection{Test it}


\begin{lstlisting}[language=bash]
# Health check
curl http://localhost:3456/health

# List models
curl http://localhost:3456/v1/models

# Chat completion
curl http://localhost:3456/v1/chat/completions \
  -H "Content-Type: application/json" \
  -d '{
    "model": "claude-opus-4",
    "messages": [{"role": "user", "content": "Hello!"}]
  }'
\end{lstlisting}


\subsubsection{With OpenClaw}


You can point OpenClaw at the proxy as a custom OpenAI-compatible endpoint:

\begin{lstlisting}[language=json]
{
  env: {
    OPENAI_API_KEY: "not-needed",
    OPENAI_BASE_URL: "http://localhost:3456/v1",
  },
  agents: {
    defaults: {
      model: { primary: "openai/claude-opus-4" },
    },
  },
}
\end{lstlisting}


\subsection{Available Models}



\begin{table}[h]
\centering
\small
\begin{tabular}{|l|l|}
\hline
Model ID & Maps To \\
\hline
\texttt{claude-opus-4} & Claude Opus 4 \\
\hline
\texttt{claude-sonnet-4} & Claude Sonnet 4 \\
\hline
\texttt{claude-haiku-4} & Claude Haiku 4 \\
\hline
\end{tabular}
\end{table}



\subsection{Auto-Start on macOS}


Create a LaunchAgent to run the proxy automatically:

\begin{lstlisting}[language=bash]
cat > ~/Library/LaunchAgents/com.claude-max-api.plist << 'EOF'
<?xml version="1.0" encoding="UTF-8"?>
<!DOCTYPE plist PUBLIC "-//Apple//DTD PLIST 1.0//EN" "http://www.apple.com/DTDs/PropertyList-1.0.dtd">
<plist version="1.0">
<dict>
  <key>Label</key>
  <string>com.claude-max-api</string>
  <key>RunAtLoad</key>
  <true/>
  <key>KeepAlive</key>
  <true/>
  <key>ProgramArguments</key>
  <array>
    <string>/usr/local/bin/node</string>
    <string>/usr/local/lib/node_modules/claude-max-api-proxy/dist/server/standalone.js</string>
  </array>
  <key>EnvironmentVariables</key>
  <dict>
    <key>PATH</key>
    <string>/usr/local/bin:/opt/homebrew/bin:~/.local/bin:/usr/bin:/bin</string>
  </dict>
</dict>
</plist>
EOF

launchctl bootstrap gui/$(id -u) ~/Library/LaunchAgents/com.claude-max-api.plist
\end{lstlisting}


\subsection{Links}


\begin{itemize}
  \item \textbf{npm:} \href{https://www.npmjs.com/package/claude-max-api-proxy}{https://www.npmjs.com/package/claude-max-api-proxy}
  \item \textbf{GitHub:} \href{https://github.com/atalovesyou/claude-max-api-proxy}{https://github.com/atalovesyou/claude-max-api-proxy}
  \item \textbf{Issues:} \href{https://github.com/atalovesyou/claude-max-api-proxy/issues}{https://github.com/atalovesyou/claude-max-api-proxy/issues}
\end{itemize}


\subsection{Notes}


\begin{itemize}
  \item This is a \textbf{community tool}, not officially supported by Anthropic or OpenClaw
  \item Requires an active Claude Max/Pro subscription with Claude Code CLI authenticated
  \item The proxy runs locally and does not send data to any third-party servers
  \item Streaming responses are fully supported
\end{itemize}


\subsection{See Also}


\begin{itemize}
  \item \href{/providers/anthropic}{Anthropic provider} - Native OpenClaw integration with Claude setup-token or API keys
  \item \href{/providers/openai}{OpenAI provider} - For OpenAI/Codex subscriptions
\end{itemize}



\clearpage

% === Source file: providers/cloudflare-ai-gateway.md ===
\section{Cloudflare AI Gateway}


Cloudflare AI Gateway sits in front of provider APIs and lets you add analytics, caching, and controls. For Anthropic, OpenClaw uses the Anthropic Messages API through your Gateway endpoint.

\begin{itemize}
  \item Provider: \texttt{cloudflare-ai-gateway}
  \item Base URL: \texttt{https://gateway.ai.cloudflare.com/v1///anthropic}
  \item Default model: \texttt{cloudflare-ai-gateway/claude-sonnet-4-5}
  \item API key: \texttt{CLOUDFLARE\_AI\_GATEWAY\_API\_KEY} (your provider API key for requests through the Gateway)
\end{itemize}


For Anthropic models, use your Anthropic API key.

\subsection{Quick start}


\begin{enumerate}
  \item Set the provider API key and Gateway details:
\end{enumerate}


\begin{lstlisting}[language=bash]
openclaw onboard --auth-choice cloudflare-ai-gateway-api-key
\end{lstlisting}


\begin{enumerate}
  \item Set a default model:
\end{enumerate}


\begin{lstlisting}[language=json]
{
  agents: {
    defaults: {
      model: { primary: "cloudflare-ai-gateway/claude-sonnet-4-5" },
    },
  },
}
\end{lstlisting}


\subsection{Non-interactive example}


\begin{lstlisting}[language=bash]
openclaw onboard --non-interactive \
  --mode local \
  --auth-choice cloudflare-ai-gateway-api-key \
  --cloudflare-ai-gateway-account-id "your-account-id" \
  --cloudflare-ai-gateway-gateway-id "your-gateway-id" \
  --cloudflare-ai-gateway-api-key "$CLOUDFLARE_AI_GATEWAY_API_KEY"
\end{lstlisting}


\subsection{Authenticated gateways}


If you enabled Gateway authentication in Cloudflare, add the \texttt{cf-aig-authorization} header (this is in addition to your provider API key).

\begin{lstlisting}[language=json]
{
  models: {
    providers: {
      "cloudflare-ai-gateway": {
        headers: {
          "cf-aig-authorization": "Bearer <cloudflare-ai-gateway-token>",
        },
      },
    },
  },
}
\end{lstlisting}


\subsection{Environment note}


If the Gateway runs as a daemon (launchd/systemd), make sure \texttt{CLOUDFLARE\_AI\_GATEWAY\_API\_KEY} is available to that process (for example, in \texttt{\textasciitilde{}/.openclaw/.env} or via \texttt{env.shellEnv}).


\clearpage

% === Source file: providers/deepgram.md ===
\section{Deepgram (Audio Transcription)}


Deepgram is a speech-to-text API. In OpenClaw it is used for **inbound audio/voice note
transcription** via \texttt{tools.media.audio}.

When enabled, OpenClaw uploads the audio file to Deepgram and injects the transcript
into the reply pipeline (\texttt{\{\{Transcript\}\}} + \texttt{[Audio]} block). This is \textbf{not streaming};
it uses the pre-recorded transcription endpoint.

Website: \href{https://deepgram.com}{https://deepgram.com}
Docs: \href{https://developers.deepgram.com}{https://developers.deepgram.com}

\subsection{Quick start}


\begin{enumerate}
  \item Set your API key:
\end{enumerate}


\begin{lstlisting}
DEEPGRAM_API_KEY=dg_...
\end{lstlisting}


\begin{enumerate}
  \item Enable the provider:
\end{enumerate}


\begin{lstlisting}[language=json]
{
  tools: {
    media: {
      audio: {
        enabled: true,
        models: [{ provider: "deepgram", model: "nova-3" }],
      },
    },
  },
}
\end{lstlisting}


\subsection{Options}


\begin{itemize}
  \item \texttt{model}: Deepgram model id (default: \texttt{nova-3})
  \item \texttt{language}: language hint (optional)
  \item \texttt{tools.media.audio.providerOptions.deepgram.detect\_language}: enable language detection (optional)
  \item \texttt{tools.media.audio.providerOptions.deepgram.punctuate}: enable punctuation (optional)
  \item \texttt{tools.media.audio.providerOptions.deepgram.smart\_format}: enable smart formatting (optional)
\end{itemize}


Example with language:

\begin{lstlisting}[language=json]
{
  tools: {
    media: {
      audio: {
        enabled: true,
        models: [{ provider: "deepgram", model: "nova-3", language: "en" }],
      },
    },
  },
}
\end{lstlisting}


Example with Deepgram options:

\begin{lstlisting}[language=json]
{
  tools: {
    media: {
      audio: {
        enabled: true,
        providerOptions: {
          deepgram: {
            detect_language: true,
            punctuate: true,
            smart_format: true,
          },
        },
        models: [{ provider: "deepgram", model: "nova-3" }],
      },
    },
  },
}
\end{lstlisting}


\subsection{Notes}


\begin{itemize}
  \item Authentication follows the standard provider auth order; \texttt{DEEPGRAM\_API\_KEY} is the simplest path.
  \item Override endpoints or headers with \texttt{tools.media.audio.baseUrl} and \texttt{tools.media.audio.headers} when using a proxy.
  \item Output follows the same audio rules as other providers (size caps, timeouts, transcript injection).
\end{itemize}



\clearpage

% === Source file: providers/github-copilot.md ===
\section{GitHub Copilot}


\subsection{What is GitHub Copilot?}


GitHub Copilot is GitHub's AI coding assistant. It provides access to Copilot
models for your GitHub account and plan. OpenClaw can use Copilot as a model
provider in two different ways.

\subsection{Two ways to use Copilot in OpenClaw}


\subsubsection{1) Built-in GitHub Copilot provider (github-copilot)}


Use the native device-login flow to obtain a GitHub token, then exchange it for
Copilot API tokens when OpenClaw runs. This is the \textbf{default} and simplest path
because it does not require VS Code.

\subsubsection{2) Copilot Proxy plugin (copilot-proxy)}


Use the \textbf{Copilot Proxy} VS Code extension as a local bridge. OpenClaw talks to
the proxy’s \texttt{/v1} endpoint and uses the model list you configure there. Choose
this when you already run Copilot Proxy in VS Code or need to route through it.
You must enable the plugin and keep the VS Code extension running.

Use GitHub Copilot as a model provider (\texttt{github-copilot}). The login command runs
the GitHub device flow, saves an auth profile, and updates your config to use that
profile.

\subsection{CLI setup}


\begin{lstlisting}[language=bash]
openclaw models auth login-github-copilot
\end{lstlisting}


You'll be prompted to visit a URL and enter a one-time code. Keep the terminal
open until it completes.

\subsubsection{Optional flags}


\begin{lstlisting}[language=bash]
openclaw models auth login-github-copilot --profile-id github-copilot:work
openclaw models auth login-github-copilot --yes
\end{lstlisting}


\subsection{Set a default model}


\begin{lstlisting}[language=bash]
openclaw models set github-copilot/gpt-4o
\end{lstlisting}


\subsubsection{Config snippet}


\begin{lstlisting}[language=json]
{
  agents: { defaults: { model: { primary: "github-copilot/gpt-4o" } } },
}
\end{lstlisting}


\subsection{Notes}


\begin{itemize}
  \item Requires an interactive TTY; run it directly in a terminal.
  \item Copilot model availability depends on your plan; if a model is rejected, try
\end{itemize}

another ID (for example \texttt{github-copilot/gpt-4.1}).
\begin{itemize}
  \item The login stores a GitHub token in the auth profile store and exchanges it for a
\end{itemize}

Copilot API token when OpenClaw runs.


\clearpage

% === Source file: providers/glm.md ===
\section{GLM models}


GLM is a \textbf{model family} (not a company) available through the Z.AI platform. In OpenClaw, GLM
models are accessed via the \texttt{zai} provider and model IDs like \texttt{zai/glm-5}.

\subsection{CLI setup}


\begin{lstlisting}[language=bash]
openclaw onboard --auth-choice zai-api-key
\end{lstlisting}


\subsection{Config snippet}


\begin{lstlisting}[language=json]
{
  env: { ZAI_API_KEY: "sk-..." },
  agents: { defaults: { model: { primary: "zai/glm-5" } } },
}
\end{lstlisting}


\subsection{Notes}


\begin{itemize}
  \item GLM versions and availability can change; check Z.AI's docs for the latest.
  \item Example model IDs include \texttt{glm-5}, \texttt{glm-4.7}, and \texttt{glm-4.6}.
  \item For provider details, see \href{/providers/zai}{/providers/zai}.
\end{itemize}



\clearpage

% === Source file: providers/huggingface.md ===
\section{Hugging Face (Inference)}


\href{https://huggingface.co/docs/inference-providers}{Hugging Face Inference Providers} offer OpenAI-compatible chat completions through a single router API. You get access to many models (DeepSeek, Llama, and more) with one token. OpenClaw uses the \textbf{OpenAI-compatible endpoint} (chat completions only); for text-to-image, embeddings, or speech use the \href{https://huggingface.co/docs/api-inference/quicktour}{HF inference clients} directly.

\begin{itemize}
  \item Provider: \texttt{huggingface}
  \item Auth: \texttt{HUGGINGFACE\_HUB\_TOKEN} or \texttt{HF\_TOKEN} (fine-grained token with \textbf{Make calls to Inference Providers})
  \item API: OpenAI-compatible (\texttt{https://router.huggingface.co/v1})
  \item Billing: Single HF token; \href{https://huggingface.co/docs/inference-providers/pricing}{pricing} follows provider rates with a free tier.
\end{itemize}


\subsection{Quick start}


\begin{enumerate}
  \item Create a fine-grained token at \href{https://huggingface.co/settings/tokens/new?ownUserPermissions=inference.serverless.write\&tokenType=fineGrained}{Hugging Face → Settings → Tokens} with the \textbf{Make calls to Inference Providers} permission.
  \item Run onboarding and choose \textbf{Hugging Face} in the provider dropdown, then enter your API key when prompted:
\end{enumerate}


\begin{lstlisting}[language=bash]
openclaw onboard --auth-choice huggingface-api-key
\end{lstlisting}


\begin{enumerate}
  \item In the \textbf{Default Hugging Face model} dropdown, pick the model you want (the list is loaded from the Inference API when you have a valid token; otherwise a built-in list is shown). Your choice is saved as the default model.
  \item You can also set or change the default model later in config:
\end{enumerate}


\begin{lstlisting}[language=json]
{
  agents: {
    defaults: {
      model: { primary: "huggingface/deepseek-ai/DeepSeek-R1" },
    },
  },
}
\end{lstlisting}


\subsection{Non-interactive example}


\begin{lstlisting}[language=bash]
openclaw onboard --non-interactive \
  --mode local \
  --auth-choice huggingface-api-key \
  --huggingface-api-key "$HF_TOKEN"
\end{lstlisting}


This will set \texttt{huggingface/deepseek-ai/DeepSeek-R1} as the default model.

\subsection{Environment note}


If the Gateway runs as a daemon (launchd/systemd), make sure \texttt{HUGGINGFACE\_HUB\_TOKEN} or \texttt{HF\_TOKEN}
is available to that process (for example, in \texttt{\textasciitilde{}/.openclaw/.env} or via
\texttt{env.shellEnv}).

\subsection{Model discovery and onboarding dropdown}


OpenClaw discovers models by calling the \textbf{Inference endpoint directly}:

\begin{lstlisting}[language=bash]
GET https://router.huggingface.co/v1/models
\end{lstlisting}


(Optional: send \texttt{Authorization: Bearer \$HUGGINGFACE\_HUB\_TOKEN} or \texttt{\$HF\_TOKEN} for the full list; some endpoints return a subset without auth.) The response is OpenAI-style \texttt{\{ "object": "list", "data": [ \{ "id": "Qwen/Qwen3-8B", "owned\_by": "Qwen", ... \}, ... ] \}}.

When you configure a Hugging Face API key (via onboarding, \texttt{HUGGINGFACE\_HUB\_TOKEN}, or \texttt{HF\_TOKEN}), OpenClaw uses this GET to discover available chat-completion models. During \textbf{interactive onboarding}, after you enter your token you see a \textbf{Default Hugging Face model} dropdown populated from that list (or the built-in catalog if the request fails). At runtime (e.g. Gateway startup), when a key is present, OpenClaw again calls \textbf{GET} \texttt{https://router.huggingface.co/v1/models} to refresh the catalog. The list is merged with a built-in catalog (for metadata like context window and cost). If the request fails or no key is set, only the built-in catalog is used.

\subsection{Model names and editable options}


\begin{itemize}
  \item \textbf{Name from API:} The model display name is \textbf{hydrated from GET /v1/models} when the API returns \texttt{name}, \texttt{title}, or \texttt{display\_name}; otherwise it is derived from the model id (e.g. \texttt{deepseek-ai/DeepSeek-R1} → “DeepSeek R1”).
  \item \textbf{Override display name:} You can set a custom label per model in config so it appears the way you want in the CLI and UI:
\end{itemize}


\begin{lstlisting}[language=json]
{
  agents: {
    defaults: {
      models: {
        "huggingface/deepseek-ai/DeepSeek-R1": { alias: "DeepSeek R1 (fast)" },
        "huggingface/deepseek-ai/DeepSeek-R1:cheapest": { alias: "DeepSeek R1 (cheap)" },
      },
    },
  },
}
\end{lstlisting}


\begin{itemize}
  \item \textbf{Provider / policy selection:} Append a suffix to the \textbf{model id} to choose how the router picks the backend:
  \item \textbf{\texttt{:fastest}} — highest throughput (router picks; provider choice is \textbf{locked} — no interactive backend picker).
  \item \textbf{\texttt{:cheapest}} — lowest cost per output token (router picks; provider choice is \textbf{locked}).
  \item \textbf{\texttt{:provider}} — force a specific backend (e.g. \texttt{:sambanova}, \texttt{:together}).
\end{itemize}


When you select \textbf{:cheapest} or \textbf{:fastest} (e.g. in the onboarding model dropdown), the provider is locked: the router decides by cost or speed and no optional “prefer specific backend” step is shown. You can add these as separate entries in \texttt{models.providers.huggingface.models} or set \texttt{model.primary} with the suffix. You can also set your default order in \href{https://hf.co/settings/inference-providers}{Inference Provider settings} (no suffix = use that order).

\begin{itemize}
  \item \textbf{Config merge:} Existing entries in \texttt{models.providers.huggingface.models} (e.g. in \texttt{models.json}) are kept when config is merged. So any custom \texttt{name}, \texttt{alias}, or model options you set there are preserved.
\end{itemize}


\subsection{Model IDs and configuration examples}


Model refs use the form \texttt{huggingface//} (Hub-style IDs). The list below is from \textbf{GET} \texttt{https://router.huggingface.co/v1/models}; your catalog may include more.

\textbf{Example IDs (from the inference endpoint):}


\begin{table}[h]
\centering
\small
\begin{tabular}{|l|l|}
\hline
Model & Ref (prefix with \texttt{huggingface/}) \\
\hline
DeepSeek R1 & \texttt{deepseek-ai/DeepSeek-R1} \\
\hline
DeepSeek V3.2 & \texttt{deepseek-ai/DeepSeek-V3.2} \\
\hline
Qwen3 8B & \texttt{Qwen/Qwen3-8B} \\
\hline
Qwen2.5 7B Instruct & \texttt{Qwen/Qwen2.5-7B-Instruct} \\
\hline
Qwen3 32B & \texttt{Qwen/Qwen3-32B} \\
\hline
Llama 3.3 70B Instruct & \texttt{meta-llama/Llama-3.3-70B-Instruct} \\
\hline
Llama 3.1 8B Instruct & \texttt{meta-llama/Llama-3.1-8B-Instruct} \\
\hline
GPT-OSS 120B & \texttt{openai/gpt-oss-120b} \\
\hline
GLM 4.7 & \texttt{zai-org/GLM-4.7} \\
\hline
Kimi K2.5 & \texttt{moonshotai/Kimi-K2.5} \\
\hline
\end{tabular}
\end{table}



You can append \texttt{:fastest}, \texttt{:cheapest}, or \texttt{:provider} (e.g. \texttt{:together}, \texttt{:sambanova}) to the model id. Set your default order in \href{https://hf.co/settings/inference-providers}{Inference Provider settings}; see \href{https://huggingface.co/docs/inference-providers}{Inference Providers} and \textbf{GET} \texttt{https://router.huggingface.co/v1/models} for the full list.

\subsubsection{Complete configuration examples}


\textbf{Primary DeepSeek R1 with Qwen fallback:}

\begin{lstlisting}[language=json]
{
  agents: {
    defaults: {
      model: {
        primary: "huggingface/deepseek-ai/DeepSeek-R1",
        fallbacks: ["huggingface/Qwen/Qwen3-8B"],
      },
      models: {
        "huggingface/deepseek-ai/DeepSeek-R1": { alias: "DeepSeek R1" },
        "huggingface/Qwen/Qwen3-8B": { alias: "Qwen3 8B" },
      },
    },
  },
}
\end{lstlisting}


\textbf{Qwen as default, with :cheapest and :fastest variants:}

\begin{lstlisting}[language=json]
{
  agents: {
    defaults: {
      model: { primary: "huggingface/Qwen/Qwen3-8B" },
      models: {
        "huggingface/Qwen/Qwen3-8B": { alias: "Qwen3 8B" },
        "huggingface/Qwen/Qwen3-8B:cheapest": { alias: "Qwen3 8B (cheapest)" },
        "huggingface/Qwen/Qwen3-8B:fastest": { alias: "Qwen3 8B (fastest)" },
      },
    },
  },
}
\end{lstlisting}


\textbf{DeepSeek + Llama + GPT-OSS with aliases:}

\begin{lstlisting}[language=json]
{
  agents: {
    defaults: {
      model: {
        primary: "huggingface/deepseek-ai/DeepSeek-V3.2",
        fallbacks: [
          "huggingface/meta-llama/Llama-3.3-70B-Instruct",
          "huggingface/openai/gpt-oss-120b",
        ],
      },
      models: {
        "huggingface/deepseek-ai/DeepSeek-V3.2": { alias: "DeepSeek V3.2" },
        "huggingface/meta-llama/Llama-3.3-70B-Instruct": { alias: "Llama 3.3 70B" },
        "huggingface/openai/gpt-oss-120b": { alias: "GPT-OSS 120B" },
      },
    },
  },
}
\end{lstlisting}


\textbf{Force a specific backend with :provider:}

\begin{lstlisting}[language=json]
{
  agents: {
    defaults: {
      model: { primary: "huggingface/deepseek-ai/DeepSeek-R1:together" },
      models: {
        "huggingface/deepseek-ai/DeepSeek-R1:together": { alias: "DeepSeek R1 (Together)" },
      },
    },
  },
}
\end{lstlisting}


\textbf{Multiple Qwen and DeepSeek models with policy suffixes:}

\begin{lstlisting}[language=json]
{
  agents: {
    defaults: {
      model: { primary: "huggingface/Qwen/Qwen2.5-7B-Instruct:cheapest" },
      models: {
        "huggingface/Qwen/Qwen2.5-7B-Instruct": { alias: "Qwen2.5 7B" },
        "huggingface/Qwen/Qwen2.5-7B-Instruct:cheapest": { alias: "Qwen2.5 7B (cheap)" },
        "huggingface/deepseek-ai/DeepSeek-R1:fastest": { alias: "DeepSeek R1 (fast)" },
        "huggingface/meta-llama/Llama-3.1-8B-Instruct": { alias: "Llama 3.1 8B" },
      },
    },
  },
}
\end{lstlisting}



\clearpage

% === Source file: providers/index.md ===
\section{Model Providers}


OpenClaw can use many LLM providers. Pick a provider, authenticate, then set the
default model as \texttt{provider/model}.

Looking for chat channel docs (WhatsApp/Telegram/Discord/Slack/Mattermost (plugin)/etc.)? See \href{/channels}{Channels}.

\subsection{Quick start}


\begin{enumerate}
  \item Authenticate with the provider (usually via \texttt{openclaw onboard}).
  \item Set the default model:
\end{enumerate}


\begin{lstlisting}[language=json]
{
  agents: { defaults: { model: { primary: "anthropic/claude-opus-4-6" } } },
}
\end{lstlisting}


\subsection{Provider docs}


\begin{itemize}
  \item \href{/providers/bedrock}{Amazon Bedrock}
  \item \href{/providers/anthropic}{Anthropic (API + Claude Code CLI)}
  \item \href{/providers/cloudflare-ai-gateway}{Cloudflare AI Gateway}
  \item \href{/providers/glm}{GLM models}
  \item \href{/providers/huggingface}{Hugging Face (Inference)}
  \item \href{/providers/kilocode}{Kilocode}
  \item \href{/providers/litellm}{LiteLLM (unified gateway)}
  \item \href{/providers/minimax}{MiniMax}
  \item \href{/providers/mistral}{Mistral}
  \item \href{/providers/moonshot}{Moonshot AI (Kimi + Kimi Coding)}
  \item \href{/providers/nvidia}{NVIDIA}
  \item \href{/providers/ollama}{Ollama (local models)}
  \item \href{/providers/openai}{OpenAI (API + Codex)}
  \item \href{/providers/opencode}{OpenCode Zen}
  \item \href{/providers/openrouter}{OpenRouter}
  \item \href{/providers/qianfan}{Qianfan}
  \item \href{/providers/qwen}{Qwen (OAuth)}
  \item \href{/providers/together}{Together AI}
  \item \href{/providers/vercel-ai-gateway}{Vercel AI Gateway}
  \item \href{/providers/venice}{Venice (Venice AI, privacy-focused)}
  \item \href{/providers/vllm}{vLLM (local models)}
  \item \href{/providers/xiaomi}{Xiaomi}
  \item \href{/providers/zai}{Z.AI}
\end{itemize}


\subsection{Transcription providers}


\begin{itemize}
  \item \href{/providers/deepgram}{Deepgram (audio transcription)}
\end{itemize}


\subsection{Community tools}


\begin{itemize}
  \item \href{/providers/claude-max-api-proxy}{Claude Max API Proxy} - Community proxy for Claude subscription credentials (verify Anthropic policy/terms before use)
\end{itemize}


For the full provider catalog (xAI, Groq, Mistral, etc.) and advanced configuration,
see \href{/concepts/model-providers}{Model providers}.


\clearpage

% === Source file: providers/kilocode.md ===
\section{Kilo Gateway}


Kilo Gateway provides a \textbf{unified API} that routes requests to many models behind a single
endpoint and API key. It is OpenAI-compatible, so most OpenAI SDKs work by switching the base URL.

\subsection{Getting an API key}


\begin{enumerate}
  \item Go to \href{https://app.kilo.ai}{app.kilo.ai}
  \item Sign in or create an account
  \item Navigate to API Keys and generate a new key
\end{enumerate}


\subsection{CLI setup}


\begin{lstlisting}[language=bash]
openclaw onboard --kilocode-api-key <key>
\end{lstlisting}


Or set the environment variable:

\begin{lstlisting}[language=bash]
export KILOCODE_API_KEY="your-api-key"
\end{lstlisting}


\subsection{Config snippet}


\begin{lstlisting}[language=json]
{
  env: { KILOCODE_API_KEY: "sk-..." },
  agents: {
    defaults: {
      model: { primary: "kilocode/anthropic/claude-opus-4.6" },
    },
  },
}
\end{lstlisting}


\subsection{Surfaced model refs}


The built-in Kilo Gateway catalog currently surfaces these model refs:

\begin{itemize}
  \item \texttt{kilocode/anthropic/claude-opus-4.6} (default)
  \item \texttt{kilocode/z-ai/glm-5:free}
  \item \texttt{kilocode/minimax/minimax-m2.5:free}
  \item \texttt{kilocode/anthropic/claude-sonnet-4.5}
  \item \texttt{kilocode/openai/gpt-5.2}
  \item \texttt{kilocode/google/gemini-3-pro-preview}
  \item \texttt{kilocode/google/gemini-3-flash-preview}
  \item \texttt{kilocode/x-ai/grok-code-fast-1}
  \item \texttt{kilocode/moonshotai/kimi-k2.5}
\end{itemize}


\subsection{Notes}


\begin{itemize}
  \item Model refs are \texttt{kilocode//} (e.g., \texttt{kilocode/anthropic/claude-opus-4.6}).
  \item Default model: \texttt{kilocode/anthropic/claude-opus-4.6}
  \item Base URL: \texttt{https://api.kilo.ai/api/gateway/}
  \item For more model/provider options, see \href{/concepts/model-providers}{/concepts/model-providers}.
  \item Kilo Gateway uses a Bearer token with your API key under the hood.
\end{itemize}



\clearpage

% === Source file: providers/litellm.md ===
\section{LiteLLM}


\href{https://litellm.ai}{LiteLLM} is an open-source LLM gateway that provides a unified API to 100+ model providers. Route OpenClaw through LiteLLM to get centralized cost tracking, logging, and the flexibility to switch backends without changing your OpenClaw config.

\subsection{Why use LiteLLM with OpenClaw?}


\begin{itemize}
  \item \textbf{Cost tracking} — See exactly what OpenClaw spends across all models
  \item \textbf{Model routing} — Switch between Claude, GPT-4, Gemini, Bedrock without config changes
  \item \textbf{Virtual keys} — Create keys with spend limits for OpenClaw
  \item \textbf{Logging} — Full request/response logs for debugging
  \item \textbf{Fallbacks} — Automatic failover if your primary provider is down
\end{itemize}


\subsection{Quick start}


\subsubsection{Via onboarding}


\begin{lstlisting}[language=bash]
openclaw onboard --auth-choice litellm-api-key
\end{lstlisting}


\subsubsection{Manual setup}


\begin{enumerate}
  \item Start LiteLLM Proxy:
\end{enumerate}


\begin{lstlisting}[language=bash]
pip install 'litellm[proxy]'
litellm --model claude-opus-4-6
\end{lstlisting}


\begin{enumerate}
  \item Point OpenClaw to LiteLLM:
\end{enumerate}


\begin{lstlisting}[language=bash]
export LITELLM_API_KEY="your-litellm-key"

openclaw
\end{lstlisting}


That's it. OpenClaw now routes through LiteLLM.

\subsection{Configuration}


\subsubsection{Environment variables}


\begin{lstlisting}[language=bash]
export LITELLM_API_KEY="sk-litellm-key"
\end{lstlisting}


\subsubsection{Config file}


\begin{lstlisting}[language=json]
{
  models: {
    providers: {
      litellm: {
        baseUrl: "http://localhost:4000",
        apiKey: "${LITELLM_API_KEY}",
        api: "openai-completions",
        models: [
          {
            id: "claude-opus-4-6",
            name: "Claude Opus 4.6",
            reasoning: true,
            input: ["text", "image"],
            contextWindow: 200000,
            maxTokens: 64000,
          },
          {
            id: "gpt-4o",
            name: "GPT-4o",
            reasoning: false,
            input: ["text", "image"],
            contextWindow: 128000,
            maxTokens: 8192,
          },
        ],
      },
    },
  },
  agents: {
    defaults: {
      model: { primary: "litellm/claude-opus-4-6" },
    },
  },
}
\end{lstlisting}


\subsection{Virtual keys}


Create a dedicated key for OpenClaw with spend limits:

\begin{lstlisting}[language=bash]
curl -X POST "http://localhost:4000/key/generate" \
  -H "Authorization: Bearer $LITELLM_MASTER_KEY" \
  -H "Content-Type: application/json" \
  -d '{
    "key_alias": "openclaw",
    "max_budget": 50.00,
    "budget_duration": "monthly"
  }'
\end{lstlisting}


Use the generated key as \texttt{LITELLM\_API\_KEY}.

\subsection{Model routing}


LiteLLM can route model requests to different backends. Configure in your LiteLLM \texttt{config.yaml}:

\begin{lstlisting}[language=yaml]
model_list:
  - model_name: claude-opus-4-6
    litellm_params:
      model: claude-opus-4-6
      api_key: os.environ/ANTHROPIC_API_KEY

  - model_name: gpt-4o
    litellm_params:
      model: gpt-4o
      api_key: os.environ/OPENAI_API_KEY
\end{lstlisting}


OpenClaw keeps requesting \texttt{claude-opus-4-6} — LiteLLM handles the routing.

\subsection{Viewing usage}


Check LiteLLM's dashboard or API:

\begin{lstlisting}[language=bash]
# Key info
curl "http://localhost:4000/key/info" \
  -H "Authorization: Bearer sk-litellm-key"

# Spend logs
curl "http://localhost:4000/spend/logs" \
  -H "Authorization: Bearer $LITELLM_MASTER_KEY"
\end{lstlisting}


\subsection{Notes}


\begin{itemize}
  \item LiteLLM runs on \texttt{http://localhost:4000} by default
  \item OpenClaw connects via the OpenAI-compatible \texttt{/v1/chat/completions} endpoint
  \item All OpenClaw features work through LiteLLM — no limitations
\end{itemize}


\subsection{See also}


\begin{itemize}
  \item \href{https://docs.litellm.ai}{LiteLLM Docs}
  \item \href{/concepts/model-providers}{Model Providers}
\end{itemize}



\clearpage

% === Source file: providers/minimax.md ===
\section{MiniMax}


MiniMax is an AI company that builds the \textbf{M2/M2.5} model family. The current
coding-focused release is \textbf{MiniMax M2.5} (December 23, 2025), built for
real-world complex tasks.

Source: \href{https://www.minimax.io/news/minimax-m25}{MiniMax M2.5 release note}

\subsection{Model overview (M2.5)}


MiniMax highlights these improvements in M2.5:

\begin{itemize}
  \item Stronger \textbf{multi-language coding} (Rust, Java, Go, C++, Kotlin, Objective-C, TS/JS).
  \item Better \textbf{web/app development} and aesthetic output quality (including native mobile).
  \item Improved \textbf{composite instruction} handling for office-style workflows, building on
\end{itemize}

interleaved thinking and integrated constraint execution.
\begin{itemize}
  \item \textbf{More concise responses} with lower token usage and faster iteration loops.
  \item Stronger \textbf{tool/agent framework} compatibility and context management (Claude Code,
\end{itemize}

Droid/Factory AI, Cline, Kilo Code, Roo Code, BlackBox).
\begin{itemize}
  \item Higher-quality \textbf{dialogue and technical writing} outputs.
\end{itemize}


\subsection{MiniMax M2.5 vs MiniMax M2.5 Highspeed}


\begin{itemize}
  \item \textbf{Speed:} \texttt{MiniMax-M2.5-highspeed} is the official fast tier in MiniMax docs.
  \item \textbf{Cost:} MiniMax pricing lists the same input cost and a higher output cost for highspeed.
  \item \textbf{Compatibility:} OpenClaw still accepts legacy \texttt{MiniMax-M2.5-Lightning} configs, but prefer
\end{itemize}

\texttt{MiniMax-M2.5-highspeed} for new setup.

\subsection{Choose a setup}


\subsubsection{MiniMax OAuth (Coding Plan) — recommended}


\textbf{Best for:} quick setup with MiniMax Coding Plan via OAuth, no API key required.

Enable the bundled OAuth plugin and authenticate:

\begin{lstlisting}[language=bash]
openclaw plugins enable minimax-portal-auth  # skip if already loaded.
openclaw gateway restart  # restart if gateway is already running
openclaw onboard --auth-choice minimax-portal
\end{lstlisting}


You will be prompted to select an endpoint:

\begin{itemize}
  \item \textbf{Global} - International users (\texttt{api.minimax.io})
  \item \textbf{CN} - Users in China (\texttt{api.minimaxi.com})
\end{itemize}


See \href{https://github.com/openclaw/openclaw/tree/main/extensions/minimax-portal-auth}{MiniMax OAuth plugin README} for details.

\subsubsection{MiniMax M2.5 (API key)}


\textbf{Best for:} hosted MiniMax with Anthropic-compatible API.

Configure via CLI:

\begin{itemize}
  \item Run \texttt{openclaw configure}
  \item Select \textbf{Model/auth}
  \item Choose \textbf{MiniMax M2.5}
\end{itemize}


\begin{lstlisting}[language=json]
{
  env: { MINIMAX_API_KEY: "sk-..." },
  agents: { defaults: { model: { primary: "minimax/MiniMax-M2.5" } } },
  models: {
    mode: "merge",
    providers: {
      minimax: {
        baseUrl: "https://api.minimax.io/anthropic",
        apiKey: "${MINIMAX_API_KEY}",
        api: "anthropic-messages",
        models: [
          {
            id: "MiniMax-M2.5",
            name: "MiniMax M2.5",
            reasoning: true,
            input: ["text"],
            cost: { input: 0.3, output: 1.2, cacheRead: 0.03, cacheWrite: 0.12 },
            contextWindow: 200000,
            maxTokens: 8192,
          },
          {
            id: "MiniMax-M2.5-highspeed",
            name: "MiniMax M2.5 Highspeed",
            reasoning: true,
            input: ["text"],
            cost: { input: 0.3, output: 1.2, cacheRead: 0.03, cacheWrite: 0.12 },
            contextWindow: 200000,
            maxTokens: 8192,
          },
        ],
      },
    },
  },
}
\end{lstlisting}


\subsubsection{MiniMax M2.5 as fallback (example)}


\textbf{Best for:} keep your strongest latest-generation model as primary, fail over to MiniMax M2.5.
Example below uses Opus as a concrete primary; swap to your preferred latest-gen primary model.

\begin{lstlisting}[language=json]
{
  env: { MINIMAX_API_KEY: "sk-..." },
  agents: {
    defaults: {
      models: {
        "anthropic/claude-opus-4-6": { alias: "primary" },
        "minimax/MiniMax-M2.5": { alias: "minimax" },
      },
      model: {
        primary: "anthropic/claude-opus-4-6",
        fallbacks: ["minimax/MiniMax-M2.5"],
      },
    },
  },
}
\end{lstlisting}


\subsubsection{Optional: Local via LM Studio (manual)}


\textbf{Best for:} local inference with LM Studio.
We have seen strong results with MiniMax M2.5 on powerful hardware (e.g. a
desktop/server) using LM Studio's local server.

Configure manually via \texttt{openclaw.json}:

\begin{lstlisting}[language=json]
{
  agents: {
    defaults: {
      model: { primary: "lmstudio/minimax-m2.5-gs32" },
      models: { "lmstudio/minimax-m2.5-gs32": { alias: "Minimax" } },
    },
  },
  models: {
    mode: "merge",
    providers: {
      lmstudio: {
        baseUrl: "http://127.0.0.1:1234/v1",
        apiKey: "lmstudio",
        api: "openai-responses",
        models: [
          {
            id: "minimax-m2.5-gs32",
            name: "MiniMax M2.5 GS32",
            reasoning: false,
            input: ["text"],
            cost: { input: 0, output: 0, cacheRead: 0, cacheWrite: 0 },
            contextWindow: 196608,
            maxTokens: 8192,
          },
        ],
      },
    },
  },
}
\end{lstlisting}


\subsection{Configure via openclaw configure}


Use the interactive config wizard to set MiniMax without editing JSON:

\begin{enumerate}
  \item Run \texttt{openclaw configure}.
  \item Select \textbf{Model/auth}.
  \item Choose \textbf{MiniMax M2.5}.
  \item Pick your default model when prompted.
\end{enumerate}


\subsection{Configuration options}


\begin{itemize}
  \item \texttt{models.providers.minimax.baseUrl}: prefer \texttt{https://api.minimax.io/anthropic} (Anthropic-compatible); \texttt{https://api.minimax.io/v1} is optional for OpenAI-compatible payloads.
  \item \texttt{models.providers.minimax.api}: prefer \texttt{anthropic-messages}; \texttt{openai-completions} is optional for OpenAI-compatible payloads.
  \item \texttt{models.providers.minimax.apiKey}: MiniMax API key (\texttt{MINIMAX\_API\_KEY}).
  \item \texttt{models.providers.minimax.models}: define \texttt{id}, \texttt{name}, \texttt{reasoning}, \texttt{contextWindow}, \texttt{maxTokens}, \texttt{cost}.
  \item \texttt{agents.defaults.models}: alias models you want in the allowlist.
  \item \texttt{models.mode}: keep \texttt{merge} if you want to add MiniMax alongside built-ins.
\end{itemize}


\subsection{Notes}


\begin{itemize}
  \item Model refs are \texttt{minimax/}.
  \item Recommended model IDs: \texttt{MiniMax-M2.5} and \texttt{MiniMax-M2.5-highspeed}.
  \item Coding Plan usage API: \texttt{https://api.minimaxi.com/v1/api/openplatform/coding\_plan/remains} (requires a coding plan key).
  \item Update pricing values in \texttt{models.json} if you need exact cost tracking.
  \item Referral link for MiniMax Coding Plan (10\% off): \href{https://platform.minimax.io/subscribe/coding-plan?code=DbXJTRClnb\&source=link}{https://platform.minimax.io/subscribe/coding-plan?code=DbXJTRClnb\&source=link}
  \item See \href{/concepts/model-providers}{/concepts/model-providers} for provider rules.
  \item Use \texttt{openclaw models list} and \texttt{openclaw models set minimax/MiniMax-M2.5} to switch.
\end{itemize}


\subsection{Troubleshooting}


\subsubsection{“Unknown model: minimax/MiniMax-M2.5”}


This usually means the \textbf{MiniMax provider isn’t configured} (no provider entry
and no MiniMax auth profile/env key found). A fix for this detection is in
\textbf{2026.1.12} (unreleased at the time of writing). Fix by:

\begin{itemize}
  \item Upgrading to \textbf{2026.1.12} (or run from source \texttt{main}), then restarting the gateway.
  \item Running \texttt{openclaw configure} and selecting \textbf{MiniMax M2.5}, or
  \item Adding the \texttt{models.providers.minimax} block manually, or
  \item Setting \texttt{MINIMAX\_API\_KEY} (or a MiniMax auth profile) so the provider can be injected.
\end{itemize}


Make sure the model id is \textbf{case‑sensitive}:

\begin{itemize}
  \item \texttt{minimax/MiniMax-M2.5}
  \item \texttt{minimax/MiniMax-M2.5-highspeed}
  \item \texttt{minimax/MiniMax-M2.5-Lightning} (legacy)
\end{itemize}


Then recheck with:

\begin{lstlisting}[language=bash]
openclaw models list
\end{lstlisting}



\clearpage

% === Source file: providers/mistral.md ===
\section{Mistral}


OpenClaw supports Mistral for both text/image model routing (\texttt{mistral/...}) and
audio transcription via Voxtral in media understanding.
Mistral can also be used for memory embeddings (\texttt{memorySearch.provider = "mistral"}).

\subsection{CLI setup}


\begin{lstlisting}[language=bash]
openclaw onboard --auth-choice mistral-api-key
# or non-interactive
openclaw onboard --mistral-api-key "$MISTRAL_API_KEY"
\end{lstlisting}


\subsection{Config snippet (LLM provider)}


\begin{lstlisting}[language=json]
{
  env: { MISTRAL_API_KEY: "sk-..." },
  agents: { defaults: { model: { primary: "mistral/mistral-large-latest" } } },
}
\end{lstlisting}


\subsection{Config snippet (audio transcription with Voxtral)}


\begin{lstlisting}[language=json]
{
  tools: {
    media: {
      audio: {
        enabled: true,
        models: [{ provider: "mistral", model: "voxtral-mini-latest" }],
      },
    },
  },
}
\end{lstlisting}


\subsection{Notes}


\begin{itemize}
  \item Mistral auth uses \texttt{MISTRAL\_API\_KEY}.
  \item Provider base URL defaults to \texttt{https://api.mistral.ai/v1}.
  \item Onboarding default model is \texttt{mistral/mistral-large-latest}.
  \item Media-understanding default audio model for Mistral is \texttt{voxtral-mini-latest}.
  \item Media transcription path uses \texttt{/v1/audio/transcriptions}.
  \item Memory embeddings path uses \texttt{/v1/embeddings} (default model: \texttt{mistral-embed}).
\end{itemize}



\clearpage

% === Source file: providers/models.md ===
\section{Model Providers}


OpenClaw can use many LLM providers. Pick one, authenticate, then set the default
model as \texttt{provider/model}.

\subsection{Quick start (two steps)}


\begin{enumerate}
  \item Authenticate with the provider (usually via \texttt{openclaw onboard}).
  \item Set the default model:
\end{enumerate}


\begin{lstlisting}[language=json]
{
  agents: { defaults: { model: { primary: "anthropic/claude-opus-4-6" } } },
}
\end{lstlisting}


\subsection{Supported providers (starter set)}


\begin{itemize}
  \item \href{/providers/openai}{OpenAI (API + Codex)}
  \item \href{/providers/anthropic}{Anthropic (API + Claude Code CLI)}
  \item \href{/providers/openrouter}{OpenRouter}
  \item \href{/providers/vercel-ai-gateway}{Vercel AI Gateway}
  \item \href{/providers/cloudflare-ai-gateway}{Cloudflare AI Gateway}
  \item \href{/providers/moonshot}{Moonshot AI (Kimi + Kimi Coding)}
  \item \href{/providers/mistral}{Mistral}
  \item \href{/providers/synthetic}{Synthetic}
  \item \href{/providers/opencode}{OpenCode Zen}
  \item \href{/providers/zai}{Z.AI}
  \item \href{/providers/glm}{GLM models}
  \item \href{/providers/minimax}{MiniMax}
  \item \href{/providers/venice}{Venice (Venice AI)}
  \item \href{/providers/bedrock}{Amazon Bedrock}
  \item \href{/providers/qianfan}{Qianfan}
\end{itemize}


For the full provider catalog (xAI, Groq, Mistral, etc.) and advanced configuration,
see \href{/concepts/model-providers}{Model providers}.


\clearpage

% === Source file: providers/moonshot.md ===
\section{Moonshot AI (Kimi)}


Moonshot provides the Kimi API with OpenAI-compatible endpoints. Configure the
provider and set the default model to \texttt{moonshot/kimi-k2.5}, or use
Kimi Coding with \texttt{kimi-coding/k2p5}.

Current Kimi K2 model IDs:


\{/\_ moonshot-kimi-k2-ids:start \_/ \&\& null\}


\begin{itemize}
  \item \texttt{kimi-k2.5}
  \item \texttt{kimi-k2-0905-preview}
  \item \texttt{kimi-k2-turbo-preview}
  \item \texttt{kimi-k2-thinking}
  \item \texttt{kimi-k2-thinking-turbo}
\end{itemize}

\{/\_ moonshot-kimi-k2-ids:end \_/ \&\& null\}

\begin{lstlisting}[language=bash]
openclaw onboard --auth-choice moonshot-api-key
\end{lstlisting}


Kimi Coding:

\begin{lstlisting}[language=bash]
openclaw onboard --auth-choice kimi-code-api-key
\end{lstlisting}


Note: Moonshot and Kimi Coding are separate providers. Keys are not interchangeable, endpoints differ, and model refs differ (Moonshot uses \texttt{moonshot/...}, Kimi Coding uses \texttt{kimi-coding/...}).

\subsection{Config snippet (Moonshot API)}


\begin{lstlisting}[language=json]
{
  env: { MOONSHOT_API_KEY: "sk-..." },
  agents: {
    defaults: {
      model: { primary: "moonshot/kimi-k2.5" },
      models: {
        // moonshot-kimi-k2-aliases:start
        "moonshot/kimi-k2.5": { alias: "Kimi K2.5" },
        "moonshot/kimi-k2-0905-preview": { alias: "Kimi K2" },
        "moonshot/kimi-k2-turbo-preview": { alias: "Kimi K2 Turbo" },
        "moonshot/kimi-k2-thinking": { alias: "Kimi K2 Thinking" },
        "moonshot/kimi-k2-thinking-turbo": { alias: "Kimi K2 Thinking Turbo" },
        // moonshot-kimi-k2-aliases:end
      },
    },
  },
  models: {
    mode: "merge",
    providers: {
      moonshot: {
        baseUrl: "https://api.moonshot.ai/v1",
        apiKey: "${MOONSHOT_API_KEY}",
        api: "openai-completions",
        models: [
          // moonshot-kimi-k2-models:start
          {
            id: "kimi-k2.5",
            name: "Kimi K2.5",
            reasoning: false,
            input: ["text"],
            cost: { input: 0, output: 0, cacheRead: 0, cacheWrite: 0 },
            contextWindow: 256000,
            maxTokens: 8192,
          },
          {
            id: "kimi-k2-0905-preview",
            name: "Kimi K2 0905 Preview",
            reasoning: false,
            input: ["text"],
            cost: { input: 0, output: 0, cacheRead: 0, cacheWrite: 0 },
            contextWindow: 256000,
            maxTokens: 8192,
          },
          {
            id: "kimi-k2-turbo-preview",
            name: "Kimi K2 Turbo",
            reasoning: false,
            input: ["text"],
            cost: { input: 0, output: 0, cacheRead: 0, cacheWrite: 0 },
            contextWindow: 256000,
            maxTokens: 8192,
          },
          {
            id: "kimi-k2-thinking",
            name: "Kimi K2 Thinking",
            reasoning: true,
            input: ["text"],
            cost: { input: 0, output: 0, cacheRead: 0, cacheWrite: 0 },
            contextWindow: 256000,
            maxTokens: 8192,
          },
          {
            id: "kimi-k2-thinking-turbo",
            name: "Kimi K2 Thinking Turbo",
            reasoning: true,
            input: ["text"],
            cost: { input: 0, output: 0, cacheRead: 0, cacheWrite: 0 },
            contextWindow: 256000,
            maxTokens: 8192,
          },
          // moonshot-kimi-k2-models:end
        ],
      },
    },
  },
}
\end{lstlisting}


\subsection{Kimi Coding}


\begin{lstlisting}[language=json]
{
  env: { KIMI_API_KEY: "sk-..." },
  agents: {
    defaults: {
      model: { primary: "kimi-coding/k2p5" },
      models: {
        "kimi-coding/k2p5": { alias: "Kimi K2.5" },
      },
    },
  },
}
\end{lstlisting}


\subsection{Notes}


\begin{itemize}
  \item Moonshot model refs use \texttt{moonshot/}. Kimi Coding model refs use \texttt{kimi-coding/}.
  \item Override pricing and context metadata in \texttt{models.providers} if needed.
  \item If Moonshot publishes different context limits for a model, adjust
\end{itemize}

\texttt{contextWindow} accordingly.
\begin{itemize}
  \item Use \texttt{https://api.moonshot.ai/v1} for the international endpoint, and \texttt{https://api.moonshot.cn/v1} for the China endpoint.
\end{itemize}


\subsection{Native thinking mode (Moonshot)}


Moonshot Kimi supports binary native thinking:

\begin{itemize}
  \item \texttt{thinking: \{ type: "enabled" \}}
  \item \texttt{thinking: \{ type: "disabled" \}}
\end{itemize}


Configure it per model via \texttt{agents.defaults.models..params}:

\begin{lstlisting}[language=json]
{
  agents: {
    defaults: {
      models: {
        "moonshot/kimi-k2.5": {
          params: {
            thinking: { type: "disabled" },
          },
        },
      },
    },
  },
}
\end{lstlisting}


OpenClaw also maps runtime \texttt{/think} levels for Moonshot:

\begin{itemize}
  \item \texttt{/think off} -> \texttt{thinking.type=disabled}
  \item any non-off thinking level -> \texttt{thinking.type=enabled}
\end{itemize}


When Moonshot thinking is enabled, \texttt{tool\_choice} must be \texttt{auto} or \texttt{none}. OpenClaw normalizes incompatible \texttt{tool\_choice} values to \texttt{auto} for compatibility.


\clearpage

% === Source file: providers/nvidia.md ===
\section{NVIDIA}


NVIDIA provides an OpenAI-compatible API at \texttt{https://integrate.api.nvidia.com/v1} for Nemotron and NeMo models. Authenticate with an API key from \href{https://catalog.ngc.nvidia.com/}{NVIDIA NGC}.

\subsection{CLI setup}


Export the key once, then run onboarding and set an NVIDIA model:

\begin{lstlisting}[language=bash]
export NVIDIA_API_KEY="nvapi-..."
openclaw onboard --auth-choice skip
openclaw models set nvidia/nvidia/llama-3.1-nemotron-70b-instruct
\end{lstlisting}


If you still pass \texttt{--token}, remember it lands in shell history and \texttt{ps} output; prefer the env var when possible.

\subsection{Config snippet}


\begin{lstlisting}[language=json]
{
  env: { NVIDIA_API_KEY: "nvapi-..." },
  models: {
    providers: {
      nvidia: {
        baseUrl: "https://integrate.api.nvidia.com/v1",
        api: "openai-completions",
      },
    },
  },
  agents: {
    defaults: {
      model: { primary: "nvidia/nvidia/llama-3.1-nemotron-70b-instruct" },
    },
  },
}
\end{lstlisting}


\subsection{Model IDs}


\begin{itemize}
  \item \texttt{nvidia/llama-3.1-nemotron-70b-instruct} (default)
  \item \texttt{meta/llama-3.3-70b-instruct}
  \item \texttt{nvidia/mistral-nemo-minitron-8b-8k-instruct}
\end{itemize}


\subsection{Notes}


\begin{itemize}
  \item OpenAI-compatible \texttt{/v1} endpoint; use an API key from NVIDIA NGC.
  \item Provider auto-enables when \texttt{NVIDIA\_API\_KEY} is set; uses static defaults (131,072-token context window, 4,096 max tokens).
\end{itemize}



\clearpage

% === Source file: providers/ollama.md ===
\section{Ollama}


Ollama is a local LLM runtime that makes it easy to run open-source models on your machine. OpenClaw integrates with Ollama's native API (\texttt{/api/chat}), supporting streaming and tool calling, and can \textbf{auto-discover tool-capable models} when you opt in with \texttt{OLLAMA\_API\_KEY} (or an auth profile) and do not define an explicit \texttt{models.providers.ollama} entry.

\textbf{Remote Ollama users}: Do not use the \texttt{/v1} OpenAI-compatible URL (\texttt{http://host:11434/v1}) with OpenClaw. This breaks tool calling and models may output raw tool JSON as plain text. Use the native Ollama API URL instead: \texttt{baseUrl: "http://host:11434"} (no \texttt{/v1}).

\subsection{Quick start}


\begin{enumerate}
  \item Install Ollama: \href{https://ollama.ai}{https://ollama.ai}
\end{enumerate}


\begin{enumerate}
  \item Pull a model:
\end{enumerate}


\begin{lstlisting}[language=bash]
ollama pull gpt-oss:20b
# or
ollama pull llama3.3
# or
ollama pull qwen2.5-coder:32b
# or
ollama pull deepseek-r1:32b
\end{lstlisting}


\begin{enumerate}
  \item Enable Ollama for OpenClaw (any value works; Ollama doesn't require a real key):
\end{enumerate}


\begin{lstlisting}[language=bash]
# Set environment variable
export OLLAMA_API_KEY="ollama-local"

# Or configure in your config file
openclaw config set models.providers.ollama.apiKey "ollama-local"
\end{lstlisting}


\begin{enumerate}
  \item Use Ollama models:
\end{enumerate}


\begin{lstlisting}[language=json]
{
  agents: {
    defaults: {
      model: { primary: "ollama/gpt-oss:20b" },
    },
  },
}
\end{lstlisting}


\subsection{Model discovery (implicit provider)}


When you set \texttt{OLLAMA\_API\_KEY} (or an auth profile) and \textbf{do not} define \texttt{models.providers.ollama}, OpenClaw discovers models from the local Ollama instance at \texttt{http://127.0.0.1:11434}:

\begin{itemize}
  \item Queries \texttt{/api/tags} and \texttt{/api/show}
  \item Keeps only models that report \texttt{tools} capability
  \item Marks \texttt{reasoning} when the model reports \texttt{thinking}
  \item Reads \texttt{contextWindow} from \texttt{model\_info[".context\_length"]} when available
  \item Sets \texttt{maxTokens} to 10× the context window
  \item Sets all costs to \texttt{0}
\end{itemize}


This avoids manual model entries while keeping the catalog aligned with Ollama's capabilities.

To see what models are available:

\begin{lstlisting}[language=bash]
ollama list
openclaw models list
\end{lstlisting}


To add a new model, simply pull it with Ollama:

\begin{lstlisting}[language=bash]
ollama pull mistral
\end{lstlisting}


The new model will be automatically discovered and available to use.

If you set \texttt{models.providers.ollama} explicitly, auto-discovery is skipped and you must define models manually (see below).

\subsection{Configuration}


\subsubsection{Basic setup (implicit discovery)}


The simplest way to enable Ollama is via environment variable:

\begin{lstlisting}[language=bash]
export OLLAMA_API_KEY="ollama-local"
\end{lstlisting}


\subsubsection{Explicit setup (manual models)}


Use explicit config when:

\begin{itemize}
  \item Ollama runs on another host/port.
  \item You want to force specific context windows or model lists.
  \item You want to include models that do not report tool support.
\end{itemize}


\begin{lstlisting}[language=json]
{
  models: {
    providers: {
      ollama: {
        baseUrl: "http://ollama-host:11434",
        apiKey: "ollama-local",
        api: "ollama",
        models: [
          {
            id: "gpt-oss:20b",
            name: "GPT-OSS 20B",
            reasoning: false,
            input: ["text"],
            cost: { input: 0, output: 0, cacheRead: 0, cacheWrite: 0 },
            contextWindow: 8192,
            maxTokens: 8192 * 10
          }
        ]
      }
    }
  }
}
\end{lstlisting}


If \texttt{OLLAMA\_API\_KEY} is set, you can omit \texttt{apiKey} in the provider entry and OpenClaw will fill it for availability checks.

\subsubsection{Custom base URL (explicit config)}


If Ollama is running on a different host or port (explicit config disables auto-discovery, so define models manually):

\begin{lstlisting}[language=json]
{
  models: {
    providers: {
      ollama: {
        apiKey: "ollama-local",
        baseUrl: "http://ollama-host:11434", // No /v1 - use native Ollama API URL
        api: "ollama", // Set explicitly to guarantee native tool-calling behavior
      },
    },
  },
}
\end{lstlisting}


Do not add \texttt{/v1} to the URL. The \texttt{/v1} path uses OpenAI-compatible mode, where tool calling is not reliable. Use the base Ollama URL without a path suffix.

\subsubsection{Model selection}


Once configured, all your Ollama models are available:

\begin{lstlisting}[language=json]
{
  agents: {
    defaults: {
      model: {
        primary: "ollama/gpt-oss:20b",
        fallbacks: ["ollama/llama3.3", "ollama/qwen2.5-coder:32b"],
      },
    },
  },
}
\end{lstlisting}


\subsection{Advanced}


\subsubsection{Reasoning models}


OpenClaw marks models as reasoning-capable when Ollama reports \texttt{thinking} in \texttt{/api/show}:

\begin{lstlisting}[language=bash]
ollama pull deepseek-r1:32b
\end{lstlisting}


\subsubsection{Model Costs}


Ollama is free and runs locally, so all model costs are set to \$0.

\subsubsection{Streaming Configuration}


OpenClaw's Ollama integration uses the \textbf{native Ollama API} (\texttt{/api/chat}) by default, which fully supports streaming and tool calling simultaneously. No special configuration is needed.

\paragraph{Legacy OpenAI-Compatible Mode}


\textbf{Tool calling is not reliable in OpenAI-compatible mode.} Use this mode only if you need OpenAI format for a proxy and do not depend on native tool calling behavior.

If you need to use the OpenAI-compatible endpoint instead (e.g., behind a proxy that only supports OpenAI format), set \texttt{api: "openai-completions"} explicitly:

\begin{lstlisting}[language=json]
{
  models: {
    providers: {
      ollama: {
        baseUrl: "http://ollama-host:11434/v1",
        api: "openai-completions",
        injectNumCtxForOpenAICompat: true, // default: true
        apiKey: "ollama-local",
        models: [...]
      }
    }
  }
}
\end{lstlisting}


This mode may not support streaming + tool calling simultaneously. You may need to disable streaming with \texttt{params: \{ streaming: false \}} in model config.

When \texttt{api: "openai-completions"} is used with Ollama, OpenClaw injects \texttt{options.num\_ctx} by default so Ollama does not silently fall back to a 4096 context window. If your proxy/upstream rejects unknown \texttt{options} fields, disable this behavior:

\begin{lstlisting}[language=json]
{
  models: {
    providers: {
      ollama: {
        baseUrl: "http://ollama-host:11434/v1",
        api: "openai-completions",
        injectNumCtxForOpenAICompat: false,
        apiKey: "ollama-local",
        models: [...]
      }
    }
  }
}
\end{lstlisting}


\subsubsection{Context windows}


For auto-discovered models, OpenClaw uses the context window reported by Ollama when available, otherwise it defaults to \texttt{8192}. You can override \texttt{contextWindow} and \texttt{maxTokens} in explicit provider config.

\subsection{Troubleshooting}


\subsubsection{Ollama not detected}


Make sure Ollama is running and that you set \texttt{OLLAMA\_API\_KEY} (or an auth profile), and that you did \textbf{not} define an explicit \texttt{models.providers.ollama} entry:

\begin{lstlisting}[language=bash]
ollama serve
\end{lstlisting}


And that the API is accessible:

\begin{lstlisting}[language=bash]
curl http://localhost:11434/api/tags
\end{lstlisting}


\subsubsection{No models available}


OpenClaw only auto-discovers models that report tool support. If your model isn't listed, either:

\begin{itemize}
  \item Pull a tool-capable model, or
  \item Define the model explicitly in \texttt{models.providers.ollama}.
\end{itemize}


To add models:

\begin{lstlisting}[language=bash]
ollama list  # See what's installed
ollama pull gpt-oss:20b  # Pull a tool-capable model
ollama pull llama3.3     # Or another model
\end{lstlisting}


\subsubsection{Connection refused}


Check that Ollama is running on the correct port:

\begin{lstlisting}[language=bash]
# Check if Ollama is running
ps aux | grep ollama

# Or restart Ollama
ollama serve
\end{lstlisting}


\subsection{See Also}


\begin{itemize}
  \item \href{/concepts/model-providers}{Model Providers} - Overview of all providers
  \item \href{/concepts/models}{Model Selection} - How to choose models
  \item \href{/gateway/configuration}{Configuration} - Full config reference
\end{itemize}



\clearpage

% === Source file: providers/openai.md ===
\section{OpenAI}


OpenAI provides developer APIs for GPT models. Codex supports \textbf{ChatGPT sign-in} for subscription
access or \textbf{API key} sign-in for usage-based access. Codex cloud requires ChatGPT sign-in.
OpenAI explicitly supports subscription OAuth usage in external tools/workflows like OpenClaw.

\subsection{Option A: OpenAI API key (OpenAI Platform)}


\textbf{Best for:} direct API access and usage-based billing.
Get your API key from the OpenAI dashboard.

\subsubsection{CLI setup}


\begin{lstlisting}[language=bash]
openclaw onboard --auth-choice openai-api-key
# or non-interactive
openclaw onboard --openai-api-key "$OPENAI_API_KEY"
\end{lstlisting}


\subsubsection{Config snippet}


\begin{lstlisting}[language=json]
{
  env: { OPENAI_API_KEY: "sk-..." },
  agents: { defaults: { model: { primary: "openai/gpt-5.2" } } },
}
\end{lstlisting}


\subsection{Option B: OpenAI Code (Codex) subscription}


\textbf{Best for:} using ChatGPT/Codex subscription access instead of an API key.
Codex cloud requires ChatGPT sign-in, while the Codex CLI supports ChatGPT or API key sign-in.

\subsubsection{CLI setup (Codex OAuth)}


\begin{lstlisting}[language=bash]
# Run Codex OAuth in the wizard
openclaw onboard --auth-choice openai-codex

# Or run OAuth directly
openclaw models auth login --provider openai-codex
\end{lstlisting}


\subsubsection{Config snippet (Codex subscription)}


\begin{lstlisting}[language=json]
{
  agents: { defaults: { model: { primary: "openai-codex/gpt-5.3-codex" } } },
}
\end{lstlisting}


\subsubsection{Transport default}


OpenClaw uses \texttt{pi-ai} for model streaming. For both \texttt{openai/*} and
\texttt{openai-codex/*}, default transport is \texttt{"auto"} (WebSocket-first, then SSE
fallback).

You can set \texttt{agents.defaults.models..params.transport}:

\begin{itemize}
  \item \texttt{"sse"}: force SSE
  \item \texttt{"websocket"}: force WebSocket
  \item \texttt{"auto"}: try WebSocket, then fall back to SSE
\end{itemize}


For \texttt{openai/*} (Responses API), OpenClaw also enables WebSocket warm-up by
default (\texttt{openaiWsWarmup: true}) when WebSocket transport is used.

Related OpenAI docs:

\begin{itemize}
  \item \href{https://platform.openai.com/docs/guides/realtime-websocket}{Realtime API with WebSocket}
  \item \href{https://platform.openai.com/docs/guides/streaming-responses}{Streaming API responses (SSE)}
\end{itemize}


\begin{lstlisting}[language=json]
{
  agents: {
    defaults: {
      model: { primary: "openai-codex/gpt-5.3-codex" },
      models: {
        "openai-codex/gpt-5.3-codex": {
          params: {
            transport: "auto",
          },
        },
      },
    },
  },
}
\end{lstlisting}


\subsubsection{OpenAI WebSocket warm-up}


OpenAI docs describe warm-up as optional. OpenClaw enables it by default for
\texttt{openai/*} to reduce first-turn latency when using WebSocket transport.

\subsubsection{Disable warm-up}


\begin{lstlisting}[language=json]
{
  agents: {
    defaults: {
      models: {
        "openai/gpt-5.2": {
          params: {
            openaiWsWarmup: false,
          },
        },
      },
    },
  },
}
\end{lstlisting}


\subsubsection{Enable warm-up explicitly}


\begin{lstlisting}[language=json]
{
  agents: {
    defaults: {
      models: {
        "openai/gpt-5.2": {
          params: {
            openaiWsWarmup: true,
          },
        },
      },
    },
  },
}
\end{lstlisting}


\subsubsection{OpenAI Responses server-side compaction}


For direct OpenAI Responses models (\texttt{openai/*} using \texttt{api: "openai-responses"} with
\texttt{baseUrl} on \texttt{api.openai.com}), OpenClaw now auto-enables OpenAI server-side
compaction payload hints:

\begin{itemize}
  \item Forces \texttt{store: true} (unless model compat sets \texttt{supportsStore: false})
  \item Injects \texttt{context\_management: [\{ type: "compaction", compact\_threshold: ... \}]}
\end{itemize}


By default, \texttt{compact\_threshold} is \texttt{70\%} of model \texttt{contextWindow} (or \texttt{80000}
when unavailable).

\subsubsection{Enable server-side compaction explicitly}


Use this when you want to force \texttt{context\_management} injection on compatible
Responses models (for example Azure OpenAI Responses):

\begin{lstlisting}[language=json]
{
  agents: {
    defaults: {
      models: {
        "azure-openai-responses/gpt-5.2": {
          params: {
            responsesServerCompaction: true,
          },
        },
      },
    },
  },
}
\end{lstlisting}


\subsubsection{Enable with a custom threshold}


\begin{lstlisting}[language=json]
{
  agents: {
    defaults: {
      models: {
        "openai/gpt-5.2": {
          params: {
            responsesServerCompaction: true,
            responsesCompactThreshold: 120000,
          },
        },
      },
    },
  },
}
\end{lstlisting}


\subsubsection{Disable server-side compaction}


\begin{lstlisting}[language=json]
{
  agents: {
    defaults: {
      models: {
        "openai/gpt-5.2": {
          params: {
            responsesServerCompaction: false,
          },
        },
      },
    },
  },
}
\end{lstlisting}


\texttt{responsesServerCompaction} only controls \texttt{context\_management} injection.
Direct OpenAI Responses models still force \texttt{store: true} unless compat sets
\texttt{supportsStore: false}.

\subsection{Notes}


\begin{itemize}
  \item Model refs always use \texttt{provider/model} (see \href{/concepts/models}{/concepts/models}).
  \item Auth details + reuse rules are in \href{/concepts/oauth}{/concepts/oauth}.
\end{itemize}



\clearpage

% === Source file: providers/opencode.md ===
\section{OpenCode Zen}


OpenCode Zen is a \textbf{curated list of models} recommended by the OpenCode team for coding agents.
It is an optional, hosted model access path that uses an API key and the \texttt{opencode} provider.
Zen is currently in beta.

\subsection{CLI setup}


\begin{lstlisting}[language=bash]
openclaw onboard --auth-choice opencode-zen
# or non-interactive
openclaw onboard --opencode-zen-api-key "$OPENCODE_API_KEY"
\end{lstlisting}


\subsection{Config snippet}


\begin{lstlisting}[language=json]
{
  env: { OPENCODE_API_KEY: "sk-..." },
  agents: { defaults: { model: { primary: "opencode/claude-opus-4-6" } } },
}
\end{lstlisting}


\subsection{Notes}


\begin{itemize}
  \item \texttt{OPENCODE\_ZEN\_API\_KEY} is also supported.
  \item You sign in to Zen, add billing details, and copy your API key.
  \item OpenCode Zen bills per request; check the OpenCode dashboard for details.
\end{itemize}



\clearpage

% === Source file: providers/openrouter.md ===
\section{OpenRouter}


OpenRouter provides a \textbf{unified API} that routes requests to many models behind a single
endpoint and API key. It is OpenAI-compatible, so most OpenAI SDKs work by switching the base URL.

\subsection{CLI setup}


\begin{lstlisting}[language=bash]
openclaw onboard --auth-choice apiKey --token-provider openrouter --token "$OPENROUTER_API_KEY"
\end{lstlisting}


\subsection{Config snippet}


\begin{lstlisting}[language=json]
{
  env: { OPENROUTER_API_KEY: "sk-or-..." },
  agents: {
    defaults: {
      model: { primary: "openrouter/anthropic/claude-sonnet-4-5" },
    },
  },
}
\end{lstlisting}


\subsection{Notes}


\begin{itemize}
  \item Model refs are \texttt{openrouter//}.
  \item For more model/provider options, see \href{/concepts/model-providers}{/concepts/model-providers}.
  \item OpenRouter uses a Bearer token with your API key under the hood.
\end{itemize}



\clearpage

% === Source file: providers/qianfan.md ===
\section{Qianfan Provider Guide}


Qianfan is Baidu's MaaS platform, provides a \textbf{unified API} that routes requests to many models behind a single
endpoint and API key. It is OpenAI-compatible, so most OpenAI SDKs work by switching the base URL.

\subsection{Prerequisites}


\begin{enumerate}
  \item A Baidu Cloud account with Qianfan API access
  \item An API key from the Qianfan console
  \item OpenClaw installed on your system
\end{enumerate}


\subsection{Getting Your API Key}


\begin{enumerate}
  \item Visit the \href{https://console.bce.baidu.com/qianfan/ais/console/apiKey}{Qianfan Console}
  \item Create a new application or select an existing one
  \item Generate an API key (format: \texttt{bce-v3/ALTAK-...})
  \item Copy the API key for use with OpenClaw
\end{enumerate}


\subsection{CLI setup}


\begin{lstlisting}[language=bash]
openclaw onboard --auth-choice qianfan-api-key
\end{lstlisting}


\subsection{Related Documentation}


\begin{itemize}
  \item \href{/gateway/configuration}{OpenClaw Configuration}
  \item \href{/concepts/model-providers}{Model Providers}
  \item \href{/concepts/agent}{Agent Setup}
  \item \href{https://cloud.baidu.com/doc/qianfan-api/s/3m7of64lb}{Qianfan API Documentation}
\end{itemize}



\clearpage

% === Source file: providers/qwen.md ===
\section{Qwen}


Qwen provides a free-tier OAuth flow for Qwen Coder and Qwen Vision models
(2,000 requests/day, subject to Qwen rate limits).

\subsection{Enable the plugin}


\begin{lstlisting}[language=bash]
openclaw plugins enable qwen-portal-auth
\end{lstlisting}


Restart the Gateway after enabling.

\subsection{Authenticate}


\begin{lstlisting}[language=bash]
openclaw models auth login --provider qwen-portal --set-default
\end{lstlisting}


This runs the Qwen device-code OAuth flow and writes a provider entry to your
\texttt{models.json} (plus a \texttt{qwen} alias for quick switching).

\subsection{Model IDs}


\begin{itemize}
  \item \texttt{qwen-portal/coder-model}
  \item \texttt{qwen-portal/vision-model}
\end{itemize}


Switch models with:

\begin{lstlisting}[language=bash]
openclaw models set qwen-portal/coder-model
\end{lstlisting}


\subsection{Reuse Qwen Code CLI login}


If you already logged in with the Qwen Code CLI, OpenClaw will sync credentials
from \texttt{\textasciitilde{}/.qwen/oauth\_creds.json} when it loads the auth store. You still need a
\texttt{models.providers.qwen-portal} entry (use the login command above to create one).

\subsection{Notes}


\begin{itemize}
  \item Tokens auto-refresh; re-run the login command if refresh fails or access is revoked.
  \item Default base URL: \texttt{https://portal.qwen.ai/v1} (override with
\end{itemize}

\texttt{models.providers.qwen-portal.baseUrl} if Qwen provides a different endpoint).
\begin{itemize}
  \item See \href{/concepts/model-providers}{Model providers} for provider-wide rules.
\end{itemize}



\clearpage

% === Source file: providers/synthetic.md ===
\section{Synthetic}


Synthetic exposes Anthropic-compatible endpoints. OpenClaw registers it as the
\texttt{synthetic} provider and uses the Anthropic Messages API.

\subsection{Quick setup}


\begin{enumerate}
  \item Set \texttt{SYNTHETIC\_API\_KEY} (or run the wizard below).
  \item Run onboarding:
\end{enumerate}


\begin{lstlisting}[language=bash]
openclaw onboard --auth-choice synthetic-api-key
\end{lstlisting}


The default model is set to:

\begin{lstlisting}
synthetic/hf:MiniMaxAI/MiniMax-M2.5
\end{lstlisting}


\subsection{Config example}


\begin{lstlisting}[language=json]
{
  env: { SYNTHETIC_API_KEY: "sk-..." },
  agents: {
    defaults: {
      model: { primary: "synthetic/hf:MiniMaxAI/MiniMax-M2.5" },
      models: { "synthetic/hf:MiniMaxAI/MiniMax-M2.5": { alias: "MiniMax M2.5" } },
    },
  },
  models: {
    mode: "merge",
    providers: {
      synthetic: {
        baseUrl: "https://api.synthetic.new/anthropic",
        apiKey: "${SYNTHETIC_API_KEY}",
        api: "anthropic-messages",
        models: [
          {
            id: "hf:MiniMaxAI/MiniMax-M2.5",
            name: "MiniMax M2.5",
            reasoning: false,
            input: ["text"],
            cost: { input: 0, output: 0, cacheRead: 0, cacheWrite: 0 },
            contextWindow: 192000,
            maxTokens: 65536,
          },
        ],
      },
    },
  },
}
\end{lstlisting}


Note: OpenClaw's Anthropic client appends \texttt{/v1} to the base URL, so use
\texttt{https://api.synthetic.new/anthropic} (not \texttt{/anthropic/v1}). If Synthetic changes
its base URL, override \texttt{models.providers.synthetic.baseUrl}.

\subsection{Model catalog}


All models below use cost \texttt{0} (input/output/cache).


\begin{table}[h]
\centering
\small
\begin{tabular}{|l|l|l|l|l|}
\hline
Model ID & Context window & Max tokens & Reasoning & Input \\
\hline
\texttt{hf:MiniMaxAI/MiniMax-M2.5} & 192000 & 65536 & false & text \\
\hline
\texttt{hf:moonshotai/Kimi-K2-Thinking} & 256000 & 8192 & true & text \\
\hline
\texttt{hf:zai-org/GLM-4.7} & 198000 & 128000 & false & text \\
\hline
\texttt{hf:deepseek-ai/DeepSeek-R1-0528} & 128000 & 8192 & false & text \\
\hline
\texttt{hf:deepseek-ai/DeepSeek-V3-0324} & 128000 & 8192 & false & text \\
\hline
\texttt{hf:deepseek-ai/DeepSeek-V3.1} & 128000 & 8192 & false & text \\
\hline
\texttt{hf:deepseek-ai/DeepSeek-V3.1-Terminus} & 128000 & 8192 & false & text \\
\hline
\texttt{hf:deepseek-ai/DeepSeek-V3.2} & 159000 & 8192 & false & text \\
\hline
\texttt{hf:meta-llama/Llama-3.3-70B-Instruct} & 128000 & 8192 & false & text \\
\hline
\texttt{hf:meta-llama/Llama-4-Maverick-17B-128E-Instruct-FP8} & 524000 & 8192 & false & text \\
\hline
\texttt{hf:moonshotai/Kimi-K2-Instruct-0905} & 256000 & 8192 & false & text \\
\hline
\texttt{hf:openai/gpt-oss-120b} & 128000 & 8192 & false & text \\
\hline
\texttt{hf:Qwen/Qwen3-235B-A22B-Instruct-2507} & 256000 & 8192 & false & text \\
\hline
\texttt{hf:Qwen/Qwen3-Coder-480B-A35B-Instruct} & 256000 & 8192 & false & text \\
\hline
\texttt{hf:Qwen/Qwen3-VL-235B-A22B-Instruct} & 250000 & 8192 & false & text + image \\
\hline
\texttt{hf:zai-org/GLM-4.5} & 128000 & 128000 & false & text \\
\hline
\texttt{hf:zai-org/GLM-4.6} & 198000 & 128000 & false & text \\
\hline
\texttt{hf:deepseek-ai/DeepSeek-V3} & 128000 & 8192 & false & text \\
\hline
\texttt{hf:Qwen/Qwen3-235B-A22B-Thinking-2507} & 256000 & 8192 & true & text \\
\hline
\end{tabular}
\end{table}



\subsection{Notes}


\begin{itemize}
  \item Model refs use \texttt{synthetic/}.
  \item If you enable a model allowlist (\texttt{agents.defaults.models}), add every model you
\end{itemize}

plan to use.
\begin{itemize}
  \item See \href{/concepts/model-providers}{Model providers} for provider rules.
\end{itemize}



\clearpage

% === Source file: providers/together.md ===
\section{Together AI}


The \href{https://together.ai}{Together AI} provides access to leading open-source models including Llama, DeepSeek, Kimi, and more through a unified API.

\begin{itemize}
  \item Provider: \texttt{together}
  \item Auth: \texttt{TOGETHER\_API\_KEY}
  \item API: OpenAI-compatible
\end{itemize}


\subsection{Quick start}


\begin{enumerate}
  \item Set the API key (recommended: store it for the Gateway):
\end{enumerate}


\begin{lstlisting}[language=bash]
openclaw onboard --auth-choice together-api-key
\end{lstlisting}


\begin{enumerate}
  \item Set a default model:
\end{enumerate}


\begin{lstlisting}[language=json]
{
  agents: {
    defaults: {
      model: { primary: "together/moonshotai/Kimi-K2.5" },
    },
  },
}
\end{lstlisting}


\subsection{Non-interactive example}


\begin{lstlisting}[language=bash]
openclaw onboard --non-interactive \
  --mode local \
  --auth-choice together-api-key \
  --together-api-key "$TOGETHER_API_KEY"
\end{lstlisting}


This will set \texttt{together/moonshotai/Kimi-K2.5} as the default model.

\subsection{Environment note}


If the Gateway runs as a daemon (launchd/systemd), make sure \texttt{TOGETHER\_API\_KEY}
is available to that process (for example, in \texttt{\textasciitilde{}/.openclaw/.env} or via
\texttt{env.shellEnv}).

\subsection{Available models}


Together AI provides access to many popular open-source models:

\begin{itemize}
  \item \textbf{GLM 4.7 Fp8} - Default model with 200K context window
  \item \textbf{Llama 3.3 70B Instruct Turbo} - Fast, efficient instruction following
  \item \textbf{Llama 4 Scout} - Vision model with image understanding
  \item \textbf{Llama 4 Maverick} - Advanced vision and reasoning
  \item \textbf{DeepSeek V3.1} - Powerful coding and reasoning model
  \item \textbf{DeepSeek R1} - Advanced reasoning model
  \item \textbf{Kimi K2 Instruct} - High-performance model with 262K context window
\end{itemize}


All models support standard chat completions and are OpenAI API compatible.


\clearpage

% === Source file: providers/venice.md ===
\section{Venice AI (Venice highlight)}


\textbf{Venice} is our highlight Venice setup for privacy-first inference with optional anonymized access to proprietary models.

Venice AI provides privacy-focused AI inference with support for uncensored models and access to major proprietary models through their anonymized proxy. All inference is private by default—no training on your data, no logging.

\subsection{Why Venice in OpenClaw}


\begin{itemize}
  \item \textbf{Private inference} for open-source models (no logging).
  \item \textbf{Uncensored models} when you need them.
  \item \textbf{Anonymized access} to proprietary models (Opus/GPT/Gemini) when quality matters.
  \item OpenAI-compatible \texttt{/v1} endpoints.
\end{itemize}


\subsection{Privacy Modes}


Venice offers two privacy levels — understanding this is key to choosing your model:


\begin{table}[h]
\centering
\small
\begin{tabular}{|l|l|l|}
\hline
Mode & Description & Models \\
\hline
\textbf{Private} & Fully private. Prompts/responses are \textbf{never stored or logged}. Ephemeral. & Llama, Qwen, DeepSeek, Venice Uncensored, etc. \\
\hline
\textbf{Anonymized} & Proxied through Venice with metadata stripped. The underlying provider (OpenAI, Anthropic) sees anonymized requests. & Claude, GPT, Gemini, Grok, Kimi, MiniMax \\
\hline
\end{tabular}
\end{table}



\subsection{Features}


\begin{itemize}
  \item \textbf{Privacy-focused}: Choose between "private" (fully private) and "anonymized" (proxied) modes
  \item \textbf{Uncensored models}: Access to models without content restrictions
  \item \textbf{Major model access}: Use Claude, GPT-5.2, Gemini, Grok via Venice's anonymized proxy
  \item \textbf{OpenAI-compatible API}: Standard \texttt{/v1} endpoints for easy integration
  \item \textbf{Streaming}: ✅ Supported on all models
  \item \textbf{Function calling}: ✅ Supported on select models (check model capabilities)
  \item \textbf{Vision}: ✅ Supported on models with vision capability
  \item \textbf{No hard rate limits}: Fair-use throttling may apply for extreme usage
\end{itemize}


\subsection{Setup}


\subsubsection{1. Get API Key}


\begin{enumerate}
  \item Sign up at \href{https://venice.ai}{venice.ai}
  \item Go to \textbf{Settings → API Keys → Create new key}
  \item Copy your API key (format: \texttt{vapi\_xxxxxxxxxxxx})
\end{enumerate}


\subsubsection{2. Configure OpenClaw}


\textbf{Option A: Environment Variable}

\begin{lstlisting}[language=bash]
export VENICE_API_KEY="vapi_xxxxxxxxxxxx"
\end{lstlisting}


\textbf{Option B: Interactive Setup (Recommended)}

\begin{lstlisting}[language=bash]
openclaw onboard --auth-choice venice-api-key
\end{lstlisting}


This will:

\begin{enumerate}
  \item Prompt for your API key (or use existing \texttt{VENICE\_API\_KEY})
  \item Show all available Venice models
  \item Let you pick your default model
  \item Configure the provider automatically
\end{enumerate}


\textbf{Option C: Non-interactive}

\begin{lstlisting}[language=bash]
openclaw onboard --non-interactive \
  --auth-choice venice-api-key \
  --venice-api-key "vapi_xxxxxxxxxxxx"
\end{lstlisting}


\subsubsection{3. Verify Setup}


\begin{lstlisting}[language=bash]
openclaw agent --model venice/llama-3.3-70b --message "Hello, are you working?"
\end{lstlisting}


\subsection{Model Selection}


After setup, OpenClaw shows all available Venice models. Pick based on your needs:

\begin{itemize}
  \item \textbf{Default model}: \texttt{venice/llama-3.3-70b} for private, balanced performance.
  \item \textbf{High-capability option}: \texttt{venice/claude-opus-45} for hard jobs.
  \item \textbf{Privacy}: Choose "private" models for fully private inference.
  \item \textbf{Capability}: Choose "anonymized" models to access Claude, GPT, Gemini via Venice's proxy.
\end{itemize}


Change your default model anytime:

\begin{lstlisting}[language=bash]
openclaw models set venice/claude-opus-45
openclaw models set venice/llama-3.3-70b
\end{lstlisting}


List all available models:

\begin{lstlisting}[language=bash]
openclaw models list | grep venice
\end{lstlisting}


\subsection{Configure via openclaw configure}


\begin{enumerate}
  \item Run \texttt{openclaw configure}
  \item Select \textbf{Model/auth}
  \item Choose \textbf{Venice AI}
\end{enumerate}


\subsection{Which Model Should I Use?}



\begin{table}[h]
\centering
\small
\begin{tabular}{|l|l|l|}
\hline
Use Case & Recommended Model & Why \\
\hline
\textbf{General chat} & \texttt{llama-3.3-70b} & Good all-around, fully private \\
\hline
\textbf{High-capability option} & \texttt{claude-opus-45} & Higher quality for hard tasks \\
\hline
\textbf{Privacy + Claude quality} & \texttt{claude-opus-45} & Best reasoning via anonymized proxy \\
\hline
\textbf{Coding} & \texttt{qwen3-coder-480b-a35b-instruct} & Code-optimized, 262k context \\
\hline
\textbf{Vision tasks} & \texttt{qwen3-vl-235b-a22b} & Best private vision model \\
\hline
\textbf{Uncensored} & \texttt{venice-uncensored} & No content restrictions \\
\hline
\textbf{Fast + cheap} & \texttt{qwen3-4b} & Lightweight, still capable \\
\hline
\textbf{Complex reasoning} & \texttt{deepseek-v3.2} & Strong reasoning, private \\
\hline
\end{tabular}
\end{table}



\subsection{Available Models (25 Total)}


\subsubsection{Private Models (15) — Fully Private, No Logging}



\begin{table}[h]
\centering
\small
\begin{tabular}{|l|l|l|l|}
\hline
Model ID & Name & Context (tokens) & Features \\
\hline
\texttt{llama-3.3-70b} & Llama 3.3 70B & 131k & General \\
\hline
\texttt{llama-3.2-3b} & Llama 3.2 3B & 131k & Fast, lightweight \\
\hline
\texttt{hermes-3-llama-3.1-405b} & Hermes 3 Llama 3.1 405B & 131k & Complex tasks \\
\hline
\texttt{qwen3-235b-a22b-thinking-2507} & Qwen3 235B Thinking & 131k & Reasoning \\
\hline
\texttt{qwen3-235b-a22b-instruct-2507} & Qwen3 235B Instruct & 131k & General \\
\hline
\texttt{qwen3-coder-480b-a35b-instruct} & Qwen3 Coder 480B & 262k & Code \\
\hline
\texttt{qwen3-next-80b} & Qwen3 Next 80B & 262k & General \\
\hline
\texttt{qwen3-vl-235b-a22b} & Qwen3 VL 235B & 262k & Vision \\
\hline
\texttt{qwen3-4b} & Venice Small (Qwen3 4B) & 32k & Fast, reasoning \\
\hline
\texttt{deepseek-v3.2} & DeepSeek V3.2 & 163k & Reasoning \\
\hline
\texttt{venice-uncensored} & Venice Uncensored & 32k & Uncensored \\
\hline
\texttt{mistral-31-24b} & Venice Medium (Mistral) & 131k & Vision \\
\hline
\texttt{google-gemma-3-27b-it} & Gemma 3 27B Instruct & 202k & Vision \\
\hline
\texttt{openai-gpt-oss-120b} & OpenAI GPT OSS 120B & 131k & General \\
\hline
\texttt{zai-org-glm-4.7} & GLM 4.7 & 202k & Reasoning, multilingual \\
\hline
\end{tabular}
\end{table}



\subsubsection{Anonymized Models (10) — Via Venice Proxy}



\begin{table}[h]
\centering
\small
\begin{tabular}{|l|l|l|l|}
\hline
Model ID & Original & Context (tokens) & Features \\
\hline
\texttt{claude-opus-45} & Claude Opus 4.5 & 202k & Reasoning, vision \\
\hline
\texttt{claude-sonnet-45} & Claude Sonnet 4.5 & 202k & Reasoning, vision \\
\hline
\texttt{openai-gpt-52} & GPT-5.2 & 262k & Reasoning \\
\hline
\texttt{openai-gpt-52-codex} & GPT-5.2 Codex & 262k & Reasoning, vision \\
\hline
\texttt{gemini-3-pro-preview} & Gemini 3 Pro & 202k & Reasoning, vision \\
\hline
\texttt{gemini-3-flash-preview} & Gemini 3 Flash & 262k & Reasoning, vision \\
\hline
\texttt{grok-41-fast} & Grok 4.1 Fast & 262k & Reasoning, vision \\
\hline
\texttt{grok-code-fast-1} & Grok Code Fast 1 & 262k & Reasoning, code \\
\hline
\texttt{kimi-k2-thinking} & Kimi K2 Thinking & 262k & Reasoning \\
\hline
\texttt{minimax-m21} & MiniMax M2.5 & 202k & Reasoning \\
\hline
\end{tabular}
\end{table}



\subsection{Model Discovery}


OpenClaw automatically discovers models from the Venice API when \texttt{VENICE\_API\_KEY} is set. If the API is unreachable, it falls back to a static catalog.

The \texttt{/models} endpoint is public (no auth needed for listing), but inference requires a valid API key.

\subsection{Streaming \& Tool Support}



\begin{table}[h]
\centering
\small
\begin{tabular}{|l|l|}
\hline
Feature & Support \\
\hline
\textbf{Streaming} & ✅ All models \\
\hline
\textbf{Function calling} & ✅ Most models (check \texttt{supportsFunctionCalling} in API) \\
\hline
\textbf{Vision/Images} & ✅ Models marked with "Vision" feature \\
\hline
\textbf{JSON mode} & ✅ Supported via \texttt{response\_format} \\
\hline
\end{tabular}
\end{table}



\subsection{Pricing}


Venice uses a credit-based system. Check \href{https://venice.ai/pricing}{venice.ai/pricing} for current rates:

\begin{itemize}
  \item \textbf{Private models}: Generally lower cost
  \item \textbf{Anonymized models}: Similar to direct API pricing + small Venice fee
\end{itemize}


\subsection{Comparison: Venice vs Direct API}



\begin{table}[h]
\centering
\small
\begin{tabular}{|l|l|l|}
\hline
Aspect & Venice (Anonymized) & Direct API \\
\hline
\textbf{Privacy} & Metadata stripped, anonymized & Your account linked \\
\hline
\textbf{Latency} & +10-50ms (proxy) & Direct \\
\hline
\textbf{Features} & Most features supported & Full features \\
\hline
\textbf{Billing} & Venice credits & Provider billing \\
\hline
\end{tabular}
\end{table}



\subsection{Usage Examples}


\begin{lstlisting}[language=bash]
# Use default private model
openclaw agent --model venice/llama-3.3-70b --message "Quick health check"

# Use Claude via Venice (anonymized)
openclaw agent --model venice/claude-opus-45 --message "Summarize this task"

# Use uncensored model
openclaw agent --model venice/venice-uncensored --message "Draft options"

# Use vision model with image
openclaw agent --model venice/qwen3-vl-235b-a22b --message "Review attached image"

# Use coding model
openclaw agent --model venice/qwen3-coder-480b-a35b-instruct --message "Refactor this function"
\end{lstlisting}


\subsection{Troubleshooting}


\subsubsection{API key not recognized}


\begin{lstlisting}[language=bash]
echo $VENICE_API_KEY
openclaw models list | grep venice
\end{lstlisting}


Ensure the key starts with \texttt{vapi\_}.

\subsubsection{Model not available}


The Venice model catalog updates dynamically. Run \texttt{openclaw models list} to see currently available models. Some models may be temporarily offline.

\subsubsection{Connection issues}


Venice API is at \texttt{https://api.venice.ai/api/v1}. Ensure your network allows HTTPS connections.

\subsection{Config file example}


\begin{lstlisting}[language=json]
{
  env: { VENICE_API_KEY: "vapi_..." },
  agents: { defaults: { model: { primary: "venice/llama-3.3-70b" } } },
  models: {
    mode: "merge",
    providers: {
      venice: {
        baseUrl: "https://api.venice.ai/api/v1",
        apiKey: "${VENICE_API_KEY}",
        api: "openai-completions",
        models: [
          {
            id: "llama-3.3-70b",
            name: "Llama 3.3 70B",
            reasoning: false,
            input: ["text"],
            cost: { input: 0, output: 0, cacheRead: 0, cacheWrite: 0 },
            contextWindow: 131072,
            maxTokens: 8192,
          },
        ],
      },
    },
  },
}
\end{lstlisting}


\subsection{Links}


\begin{itemize}
  \item \href{https://venice.ai}{Venice AI}
  \item \href{https://docs.venice.ai}{API Documentation}
  \item \href{https://venice.ai/pricing}{Pricing}
  \item \href{https://status.venice.ai}{Status}
\end{itemize}



\clearpage

% === Source file: providers/vercel-ai-gateway.md ===
\section{Vercel AI Gateway}


The \href{https://vercel.com/ai-gateway}{Vercel AI Gateway} provides a unified API to access hundreds of models through a single endpoint.

\begin{itemize}
  \item Provider: \texttt{vercel-ai-gateway}
  \item Auth: \texttt{AI\_GATEWAY\_API\_KEY}
  \item API: Anthropic Messages compatible
\end{itemize}


\subsection{Quick start}


\begin{enumerate}
  \item Set the API key (recommended: store it for the Gateway):
\end{enumerate}


\begin{lstlisting}[language=bash]
openclaw onboard --auth-choice ai-gateway-api-key
\end{lstlisting}


\begin{enumerate}
  \item Set a default model:
\end{enumerate}


\begin{lstlisting}[language=json]
{
  agents: {
    defaults: {
      model: { primary: "vercel-ai-gateway/anthropic/claude-opus-4.6" },
    },
  },
}
\end{lstlisting}


\subsection{Non-interactive example}


\begin{lstlisting}[language=bash]
openclaw onboard --non-interactive \
  --mode local \
  --auth-choice ai-gateway-api-key \
  --ai-gateway-api-key "$AI_GATEWAY_API_KEY"
\end{lstlisting}


\subsection{Environment note}


If the Gateway runs as a daemon (launchd/systemd), make sure \texttt{AI\_GATEWAY\_API\_KEY}
is available to that process (for example, in \texttt{\textasciitilde{}/.openclaw/.env} or via
\texttt{env.shellEnv}).

\subsection{Model ID shorthand}


OpenClaw accepts Vercel Claude shorthand model refs and normalizes them at
runtime:

\begin{itemize}
  \item \texttt{vercel-ai-gateway/claude-opus-4.6} -> \texttt{vercel-ai-gateway/anthropic/claude-opus-4.6}
  \item \texttt{vercel-ai-gateway/opus-4.6} -> \texttt{vercel-ai-gateway/anthropic/claude-opus-4-6}
\end{itemize}



\clearpage

% === Source file: providers/vllm.md ===
\section{vLLM}


vLLM can serve open-source (and some custom) models via an \textbf{OpenAI-compatible} HTTP API. OpenClaw can connect to vLLM using the \texttt{openai-completions} API.

OpenClaw can also \textbf{auto-discover} available models from vLLM when you opt in with \texttt{VLLM\_API\_KEY} (any value works if your server doesn’t enforce auth) and you do not define an explicit \texttt{models.providers.vllm} entry.

\subsection{Quick start}


\begin{enumerate}
  \item Start vLLM with an OpenAI-compatible server.
\end{enumerate}


Your base URL should expose \texttt{/v1} endpoints (e.g. \texttt{/v1/models}, \texttt{/v1/chat/completions}). vLLM commonly runs on:

\begin{itemize}
  \item \texttt{http://127.0.0.1:8000/v1}
\end{itemize}


\begin{enumerate}
  \item Opt in (any value works if no auth is configured):
\end{enumerate}


\begin{lstlisting}[language=bash]
export VLLM_API_KEY="vllm-local"
\end{lstlisting}


\begin{enumerate}
  \item Select a model (replace with one of your vLLM model IDs):
\end{enumerate}


\begin{lstlisting}[language=json]
{
  agents: {
    defaults: {
      model: { primary: "vllm/your-model-id" },
    },
  },
}
\end{lstlisting}


\subsection{Model discovery (implicit provider)}


When \texttt{VLLM\_API\_KEY} is set (or an auth profile exists) and you \textbf{do not} define \texttt{models.providers.vllm}, OpenClaw will query:

\begin{itemize}
  \item \texttt{GET http://127.0.0.1:8000/v1/models}
\end{itemize}


…and convert the returned IDs into model entries.

If you set \texttt{models.providers.vllm} explicitly, auto-discovery is skipped and you must define models manually.

\subsection{Explicit configuration (manual models)}


Use explicit config when:

\begin{itemize}
  \item vLLM runs on a different host/port.
  \item You want to pin \texttt{contextWindow}/\texttt{maxTokens} values.
  \item Your server requires a real API key (or you want to control headers).
\end{itemize}


\begin{lstlisting}[language=json]
{
  models: {
    providers: {
      vllm: {
        baseUrl: "http://127.0.0.1:8000/v1",
        apiKey: "${VLLM_API_KEY}",
        api: "openai-completions",
        models: [
          {
            id: "your-model-id",
            name: "Local vLLM Model",
            reasoning: false,
            input: ["text"],
            cost: { input: 0, output: 0, cacheRead: 0, cacheWrite: 0 },
            contextWindow: 128000,
            maxTokens: 8192,
          },
        ],
      },
    },
  },
}
\end{lstlisting}


\subsection{Troubleshooting}


\begin{itemize}
  \item Check the server is reachable:
\end{itemize}


\begin{lstlisting}[language=bash]
curl http://127.0.0.1:8000/v1/models
\end{lstlisting}


\begin{itemize}
  \item If requests fail with auth errors, set a real \texttt{VLLM\_API\_KEY} that matches your server configuration, or configure the provider explicitly under \texttt{models.providers.vllm}.
\end{itemize}



\clearpage

% === Source file: providers/xiaomi.md ===
\section{Xiaomi MiMo}


Xiaomi MiMo is the API platform for \textbf{MiMo} models. It provides REST APIs compatible with
OpenAI and Anthropic formats and uses API keys for authentication. Create your API key in
the \href{https://platform.xiaomimimo.com/\#/console/api-keys}{Xiaomi MiMo console}. OpenClaw uses
the \texttt{xiaomi} provider with a Xiaomi MiMo API key.

\subsection{Model overview}


\begin{itemize}
  \item \textbf{mimo-v2-flash}: 262144-token context window, Anthropic Messages API compatible.
  \item Base URL: \texttt{https://api.xiaomimimo.com/anthropic}
  \item Authorization: \texttt{Bearer \$XIAOMI\_API\_KEY}
\end{itemize}


\subsection{CLI setup}


\begin{lstlisting}[language=bash]
openclaw onboard --auth-choice xiaomi-api-key
# or non-interactive
openclaw onboard --auth-choice xiaomi-api-key --xiaomi-api-key "$XIAOMI_API_KEY"
\end{lstlisting}


\subsection{Config snippet}


\begin{lstlisting}[language=json]
{
  env: { XIAOMI_API_KEY: "your-key" },
  agents: { defaults: { model: { primary: "xiaomi/mimo-v2-flash" } } },
  models: {
    mode: "merge",
    providers: {
      xiaomi: {
        baseUrl: "https://api.xiaomimimo.com/anthropic",
        api: "anthropic-messages",
        apiKey: "XIAOMI_API_KEY",
        models: [
          {
            id: "mimo-v2-flash",
            name: "Xiaomi MiMo V2 Flash",
            reasoning: false,
            input: ["text"],
            cost: { input: 0, output: 0, cacheRead: 0, cacheWrite: 0 },
            contextWindow: 262144,
            maxTokens: 8192,
          },
        ],
      },
    },
  },
}
\end{lstlisting}


\subsection{Notes}


\begin{itemize}
  \item Model ref: \texttt{xiaomi/mimo-v2-flash}.
  \item The provider is injected automatically when \texttt{XIAOMI\_API\_KEY} is set (or an auth profile exists).
  \item See \href{/concepts/model-providers}{/concepts/model-providers} for provider rules.
\end{itemize}



\clearpage

% === Source file: providers/zai.md ===
\section{Z.AI}


Z.AI is the API platform for \textbf{GLM} models. It provides REST APIs for GLM and uses API keys
for authentication. Create your API key in the Z.AI console. OpenClaw uses the \texttt{zai} provider
with a Z.AI API key.

\subsection{CLI setup}


\begin{lstlisting}[language=bash]
openclaw onboard --auth-choice zai-api-key
# or non-interactive
openclaw onboard --zai-api-key "$ZAI_API_KEY"
\end{lstlisting}


\subsection{Config snippet}


\begin{lstlisting}[language=json]
{
  env: { ZAI_API_KEY: "sk-..." },
  agents: { defaults: { model: { primary: "zai/glm-5" } } },
}
\end{lstlisting}


\subsection{Notes}


\begin{itemize}
  \item GLM models are available as \texttt{zai/} (example: \texttt{zai/glm-5}).
  \item \texttt{tool\_stream} is enabled by default for Z.AI tool-call streaming. Set
\end{itemize}

\texttt{agents.defaults.models["zai/"].params.tool\_stream} to \texttt{false} to disable it.
\begin{itemize}
  \item See \href{/providers/glm}{/providers/glm} for the model family overview.
  \item Z.AI uses Bearer auth with your API key.
\end{itemize}



\clearpage
